\subsection{Cleaning the ordered \toughpair}
\label{sec:cleaningordered}

The cleaning of ordered \toughpair s is done in two steps:
we first show that an ordered \toughpair\ can be identified to a so-called \emph{semi-cleaned ordered \toughpair} (see \Cref{sec:semicleaning}) and show thereafter how to further identify a semi-cleaned ordered \toughpair\ to a hard matching pattern (see \Cref{sec:semicleantoclean}).

\subsubsection{Semi-cleaning}
\label{sec:semicleaning}

Recall the definition of a semi-cleaned ordered $t$-\toughpair\ from Section~\ref{sec:cleaning}.
The goal of this section is to prove Lemma~\ref{lem:semi-cleaning} restated below.

\semicleaning*

To prove Lemma~\ref{lem:semi-cleaning}, we first define a type of a vertex in $D$ with respect to the ordered $t$-\toughpair\ $(A,B)$. This is based on the neighbourhoods of a vertex in the head sets and tail sets of $A$ and $B$. We then define a notion of a bad type. The vertices having a bad type hinder the semi-cleaning procedure. To overcome this,
we show that if a digraph $D$ contains an ordered $t^2$-\toughpair\ then it also contains an ordered $t$-\toughpair\ $(A',B')$ such that there are no bad vertices with respect to $(A',B')$. We then show that given such an ordered \toughpair\ $(A',B')$ we can identify the remaining vertices to one of the endpoints of $A' \cup B'$.

\paragraph*{Bad vertices with respect to the ordered \toughpair\ $(A,B)$.}
Given an ordered \toughpair\ $(A,B)$ in a digraph $D$,
we say that a vertex $v \in V(D) \setminus V(A\cup B)$ is {\em bad} with respect to $(A,B)$ if $v \in N^+_{D^{\star}}(\tail(A)) \cap N^-_{D^{\star}}(\head(A)) \cap N^+_{D^{\star}}(\tail(B)) \cap N^-_{D^{\star}}(\head(B))$. 

\begin{lemma}[Eliminating bad vertices]\label{lem:badvertices}
Let $D$ be a
digraph such that $D$ contains an ordered $t^2$-\toughpair. Then $D$ contains an ordered $t$-\toughpair\ $(A',B')$ such that there is no bad vertex in $V(D) \setminus V(A' \cup B')$ with respect to $(A',B')$.
\end{lemma}

\begin{proof}
Let $(A,B)$ be an ordered $t^2$-\toughpair\ in $D$ and
let $t'=t^2$.
Let $A=((w_1,x_1), \ldots, (w_{t'},x_{t'}))$ and $B=((y_1,z_1), \ldots, (y_{t'},z_{t'}))$.
For each vertex $v \in V(H) \setminus V(A \cup B)$ which is bad with respect to $(A,B)$, let $W(v)$ be the largest index in $[t']$ such that $(a_{W(v)},v) \in E(D)$, $X(v)$ be the smallest index in $[t']$ such that $(v,b_{X(v)}) \in E(D)$, $Y(v)$ be the largest index in $[t']$ such that $(c_{Y(v)},v) \in E(D)$, and $Z(v)$ be the smallest index in $[t']$ such that $(v,d_{Z(v)}) \in E(D)$.
For every $p,q \in [\sqrt{t'}]$, let $V_{p,q}$ be the set of bad vertices $v$ with respect to $(A,B)$ such that
$W(v),X(v) \in \{(p-1)\sqrt{t'}+1, \ldots, p\sqrt{t'}\}$, and $Y(v),Z(v) \in \{(q-1)\sqrt{t'}, \ldots, q\sqrt{t'}\}$.

If there exists $p,q \in [\sqrt{t'}]$ such that $V_{pq}=\emptyset$, then let $A'=((w_{(p-1)\sqrt{t'}+1},x_{(p-1)\sqrt{t'}}+1), \ldots, (x_{p\sqrt{t'}}, x_{p\sqrt{t'}}))$ and 
$B'=((c_{(q-1)\sqrt{t'}+1},z_{(q-1)\sqrt{t'}+1}),\ldots, (c_{q\sqrt{t'}},z_{q\sqrt{t'}}))$. Then from the definition of $V_{p,q}$, $(A',B')$ is an ordered $\sqrt{t'}$-\toughpair\ such that there are no bad vertices with respect to it.

Without loss of generality, assume that for each $p,q \in [\sqrt{t'}]$, $V_{pq}\neq \emptyset$. Let $v_{p,q}$ denote an arbitrarily fixed vertex in $V_{pq}$. 
Let $A'=((w_{W(v_{2,1}}, v_{2,1}), \ldots, (w_{W(v_{\sqrt{t'},1})}, v_{\sqrt{t'},1}))$ and
 $B'=((y_{Y(v_{1,2})}, v_{1,2}),\ldots, (y_{Y(v_{\sqrt{t'},1})},v_{1,\sqrt{t'}} ))$. We will now show that $(A',B')$ is an ordered \toughpair\ in $D$. Moreover, we will show that there is no $(\tail(B'),\head(A'))$-path in $D$, thereby concluding that there is no bad vertex with respect to $(A',B')$. Note that the edges of $A',B'$ are minimal edges because they belong to $D$ which is a reachability minimal digraph.

\paragraph*{Weak independence of $A'$ (resp.~$B'$).}  We show that  there is no $(\head(A'),\tail(A'))$ path (resp.~no $(\head(B'),\tail(B'))$ path) in $D$. Recall that every vertex in $\head(A')$ is a bad vertex with respect to $(A,B)$.
Thus, each vertex of $\tail(A')$ is an out-neighbour of some vertex of $\tail(B')$. 
Thus, a $(\head(A'),\tail(A'))$-path in $D$ implies a $(\tail(B),\tail(A'))$-path in $D$. Since $\tail(A') \subseteq \tail(A)$, this contradicts that every edge of $A$ is strongly independent with every edge of $B$. 
The other case is symmetric.

\paragraph*{Weak independence between $A'$ and $B'$.} We show that a pair of edges containing an edge of $A'$ and an edge of $B'$ is weakly independent.
To see that there is no $(\head(B'),\tail(A'))$-path, 
recall that every vertex of $\head(B')$ is a bad vertex with respect to $(A,B)$.
Thus, each vertex of $\head(B')$ is an out-neighbour of some vertex of $\tail(B)$. Thus, a $(\head(B'),\tail(A'))$-path in $D$ implies a $(\tail(B),\head(A'))$-path in $D$. Since $\tail(A') \subseteq \tail(A)$, this is a contradiction. One can symmetrically prove that there is no
$(\head(A'),\tail(B'))$-path.

\paragraph*{Strong independence of $A'$ and $B'$.} We prove the following three statements.

{\it There is no $(\tail(A'),\tail(B'))$-path and no $(\tail(B'),\tail(A'))$-path.} Since $\tail(A') \subseteq \tail(A)$ and $\tail(B') \subseteq \tail(B)$, and $(A,B)$ is a  \toughpair, we conclude that there is no $(\tail(A'),\allowbreak
\tail(B'))$-path and no $(\tail(B'),\tail(A'))$-path.

{\it There is no $(\head(A'),\head(B'))$-path.} For the sake of contradiction, suppose there exists $v_{i,1} \in \head(B')$ and $v_{1,j} \in \head(B')$, $i,j \neq 1$, such that there exists a $(v_{i,1}, v_{1,j})$-path in $D$. By definition $w(v_{i,1}) \in \{2\sqrt{t'}+1, \ldots, t'\}$ and $X(v_{1,j}) \in \{1, \ldots, \sqrt{t'}\}$. Also, $(w_{W(v_{i,1})}, v_{i,1}) \in E(D)$ and $(v_{1,j}, x_{X(v_{1,j})}) \in E(D)$. Thus, a $(v_{i,1}, v_{1,j})$-path in $D$ implies a $(w_{W(v_{i,1})}, y_{Y(v_{1,j})})$-path in $D$. Since $X(v_{1,j}) < W(v_{i,1})$, this is a contradiction.

{\it There is no $(\head(B'),\head(A'))$-path.} For the sake of contradiction, suppose there exists $v_{i,1} \in \head(A')$ and $v_{1,j} \in \head(B')$, $i,j \neq 1$, such that there exists a $(v_{1,j}, v_{i,1})$-path in $D$. By definition $Y(v_{1,j}) \in \{2\sqrt{t'}+1, \ldots, t'\}$ and $Z(v_{i,1}) \in \{1, \ldots, \sqrt{t'}\}$. Also, $(y_{Y(v_{1,j})}, v_{1,j}) \in E(D)$ and $(v_{i,1}, z_{Z(v_{i,1})}) \in E(D)$. Thus, a $(v_{1,j}, v_{i,1})$-path in $D$ implies a $(y_{Y(v_{1,j})}, z_{Z(v_{i,1})})$-path in $D$. Since $Z(v_{i,1}) < Y(v_{1,j})$, this is a contradiction.

\paragraph*{Ordered condition on $A'$ (resp.~$B'$).}
Consider vertices $v_{i,1}$ and $v_{j,1}$ such that $i \leq j$. Suppose for the sake of contradiction that there exists a $(v_{j,1},v_{i,1})$-path in $D$. Recall $(w_{W(v_{j,1})}, v_{j,1}) \in E(D)$ and $(v_{i,1},x_{X(v_{i,1})}) \in E(D)$. Thus, there exists a $(w_{W(v_{j,1})}, y_{Y(v_{i,1})})$-path in $D$. Since $i \leq j$, $W(v_{j,1}) \leq X(v_{i,1})$. This is a contradiction. The other case can be proved symmetrically.

\paragraph*{There is no $(\tail(B'),\head(A'))$-path.} For the sake of contradiction, let $v_{i,1} \in \head(A')$, $i\neq 1$, and $y_j \in \tail(B')$ such that there exists a $(y_j, v_{i,1})$-path in $D$. Since $y_j \in \tail(B')$, $j \in \{2\sqrt{t'} +1, \ldots, t'\}$. Also $Z(v_{i,1}) \in \{1, \ldots, \sqrt{t}\}$. Since $(v_{i,1}, z_{Z(v_{i,1})}) \in E(D)$, we conclude that there is a $(y_j, z_{Z(v_{i,1})})$-path. This is a contradiction which concludes the proof.
\end{proof}

\begin{lemma}\label{lem:semin-cleaning-wo-bad}
Let $D$ be a directed 
graph and let $(A,B)$ be an ordered $t$-\toughpair\ in $D$ such that there are no bad vertices with respect to $(A,B)$ in $D$. Then one can obtain a semi-cleaned ordered $t$-\toughpair\ from $D$ by identification.
\end{lemma}

The remainder of this section is devoted to the proof of Lemma~\ref{lem:semin-cleaning-wo-bad}. Note that Lemma~\ref{lem:semin-cleaning-wo-bad}, together with Lemma~\ref{lem:badvertices}, proves Lemma~\ref{lem:semi-cleaning}.

Recall that $A,B$ are ordered sets. Let $\tail(A)=W=(w_1, \ldots,w_t)$, $\head(A)=X=(x_1, \ldots,x_t)$, $\tail(B)=Y=(y_1, \ldots, y_t)$ and $\head(B)=Z=(z_1, \ldots, z_t)$.
We will now describe a procedure of identifying vertices in $V(D) \setminus V(A \cup B)$ onto the vertices of $V(A \cup B)$, as long as there is a vertex that has both an in-neighbour and an out-neighbour in $V(A \cup B) = W \cup X \cup Y \cup Z$ in the graph $D^*$, while maintaining the invariants
that after the identification $(A,B)$ remains an ordered \toughpair\ 
in the new graph, no vertex in the new graph is bad with respect to $(A \cup B)$ and the new graph is acyclic. We will apply these identification rules exhaustively, in order.

\begin{invariants*}\label{inv:ordered-cleaning}
If $\widehat{D}$ is the graph obtained after identification, then $(A,B)$ is an ordered \toughpair\ in $\widehat{D}$, and
there are no bad vertices with respect to $(A,B)$ in $\widehat{D}$.
\end{invariants*}

\begin{identification rule}\label{ir:A}
If there exists $v \in V(D) \setminus (W \cup X \cup Y \cup Z)$ such that $N_{D^{\star}}^+(v) \cap W \neq \emptyset$ and $N_{D^{\star}}^-(v) \cap W \neq \emptyset$, then let $i \in [t]$ be the largest index such that $(w_i,v) \in E(D^{\star})$. Identify $v$ onto $w_i$.
\end{identification rule}

\begin{lemma}\label{lem:irA-safe}
Let $\widehat{D}$ be the graph obtained after the application of Identification Rule~\ref{ir:A}. Then, the invariants are satisfied.
\end{lemma}

\begin{proof}
Since  $N_{D^{\star}}^+(v) \cap W \neq \emptyset$ and $N_{D^{\star}}^-(v) \cap W \neq \emptyset$, 
there is a $(v,W)$-path, say $P_{\out}$, and a $(W,v)$-path say $P_{\inp}$ in $D$.\\

\noindent{\bf Weak independence of $A$ in $\widehat{D}$.} For the sake of contradiction, say there is a $(X,w_i)$-path in $\widehat{D}$. Then there exists a $(X,v)$-path in $D$. This, together with the $(v,W)$-path $P_{\out}$, implies an $(X,W)$-path in $D$, which is a contradiction to the weak independence of the edges of $A$ in $D$.\\

\noindent{\bf Weak independence between $A,B$ in $\widehat{D}$.} For the sake of contradiction, say there is a $(Z,w_i)$-path in $\widehat{D}$. 
Then there exists a $(Z,v)$-path in $D$. This, together with $(v,W)$-path $P_{\out}$, implies an $(Z,W)$-path in $D$, which is a contradiction to the weak independence between the edges of $A, B$ in $D$.\\

\noindent{\bf Strong independence between $A,B$ in $\widehat{D}$.} For the sake of contradiction, 
say there is a $(w_i,Y)$-path in $\widehat{D}$. 
Then there exists a $(v,Y)$-path in $D$. This, together with the $(W,v$-path $P_{\inp}$ in $D$, implies an $(W,Y)$-path in $D$, which is a contradiction to the strong independence between the edges of $A, B$ in $D$.

For the sake of contradiction, 
say there is a $(Y,w_i)$-path in $\widehat{D}$. 
Then there exists a $(Y,v)$-path in $D$. This, together with the $(v,W)$-path $P_{\out}$, implies an $(Y,W)$-path in $D$, which is a contradiction to the strong independence between the edges of $A, B$ in $D$.

\noindent{\bf Orderedness.} Fix $j <i$. Suppose, for the sake of contradiction, that there is a $(w_i,w_j)$-path in $\widehat{D}$. Then, there is a $(v,w_j)$-path in $D$. This, together with the $(w_i,v)$-path in $D$, implies a $(w_i,w_j)$-path in $D$, contradicting the ordered-ness condition on the edges of $A$.

For the sake of contradiction, for $j <i$, say there is a $(w_i,x_j)$-path in $\widehat{D}$. Then there is a $(v,x_j)$-path in $D$, which together with the $(w_i,v)$-path in $D$ contradicts the ordered-ness condition on the edges of $A$.\\

\noindent{\bf No bad vertices.} For the sake of contradiction, say there exists a bad vertex $u$ with respect to $(A,B)$ in $\widehat{D}$, that is, $N_{\widehat{D}^{\star}}^-(u) \cap W, N_{\widehat{D}^{\star}}^+(u) \cap X, N_{\widehat{D}^{\star}}^-(u) \cap Y, N_{\widehat{D}^{\star}}^+(u) \cap Z \neq \emptyset$. 
Then, considering that $\widehat{D}$ is obtained from $D$ by only identifying $v$ onto $w_i$, we get that
$N_{D^{\star}}^+(u) \cap X, N_{D^{\star}}^-(u) \cap Y, N_{D^{\star}}^+(u) \cap Z \neq \emptyset$ and $(v,u) \in E(D^{\star})$.  The $(v,u)$-path, together with the $(W,v)$-path $P_{\in}$, implies a $(W,u)$-path in $D$, that is $N_{D^{\star}}^-(u) \cap W \neq \emptyset$. That is, $u$ is a bad vertex with respect to $(A,B)$ in $D$, which is a contradiction.
\end{proof}

Identification Rule~\ref{ir:B} is the analogue of Identification Rule~\ref{ir:A}.

\begin{identification rule}\label{ir:B}
If there exists $v \in V(D) \setminus (W \cup X \cup Y \cup Z)$ such that $N_{D^{\star}}^+(v) \cap Y \neq \emptyset$ and $N_{D^{\star}}^-(v) \cap Y \neq \emptyset$, then let $i \in [t]$ be the largest index such that $(y_i,v) \in E(D^{\star})$. Identify $v$ onto $y_i$.
\end{identification rule}

Symmetrically to \Cref{lem:irA-safe}, the following holds.

\begin{lemma}\label{lem:irB-safe}
Let $\widehat{D}$ be the graph obtained after the application of Identification Rule~\ref{ir:B}. Then, the invariants are satisfied.
\end{lemma}

\begin{identification rule}\label{ir:C}
If there exists $v \in V(D) \setminus (W \cup X \cup Y \cup Z)$ such that $N_{D^{\star}}^-(v) \cap X \neq \emptyset$ and  $N_{D^{\star}}^+(v) \cap X \neq \emptyset$,
then let $i \in [t]$ be the largest index such that either $(w_i,v) \in E(D^{\star})$ or $(x_i,v) \in E(D^{\star})$. 
Identify $v$ onto $x_i$.
\end{identification rule}

\begin{lemma}\label{lem:irC-safe}
Let $\widehat{D}$ be the graph obtained after the application of Identification Rule~\ref{ir:C}. Then, the invariants are satisfied.
\end{lemma}

\begin{proof}
Since $N_{D^{\star}}^-(v) \cap X \neq \emptyset$ and $N_{D^{\star}}^+(v) \cap X \neq \emptyset$, 
here is a $(X,v)$-path, say $P_{\inp}$, and a $(v,X)$-path say $P_{\out}$ in $D$.\\

\noindent{\bf Weak independence of $A$ in $\widehat{D}$.} For the sake of contradiction, say there is a $(x_i,w_j)$-path in $\widehat{D}$. Then there exists a $(v,w_j)$-path, say $P$, in $D$. 
From the choice of $i$ and since $N_{D^{\star}}^-(v) \cap X \neq \emptyset$, there exists $i' \leq i$, such that there is a $(x_{i'},v)$-path in $D$. This implies a $(x_{i'},w_j)$-path in $D$. Since $i' <i$ and $A$ is weakly independent in $D$, $j < i$. 

From the choice of $i$, either $(w_i,v) \in E(D^*)$, in which case the existence of $P$ implies a $(w_i,w_j)$-path, $j<i$, in $D$, contradicting the weak independence of $A$ in $D$. Or, $(x_i,v) \in E(D^{\star})$, in which case $P$ implies a $(x_i,w_j)$-path, $i <i$ in $D$, again contradicting the weak independence of $A$ in $D$.\\

\noindent{\bf Weak independence between $A,B$ in $\widehat{D}$.} For the sake of contradiction, say there is a $(x_i,Y)$-path in $\widehat{D}$. 
Then there exists a $(v,Y)$-path in $D$. 
This, together with $(X,v)$-path $P_{\inp}$, implies an $(X,Y)$-path in $D$, which is a contradiction to the weak independence between the edges of $A, B$ in $D$.\\

\noindent{\bf Strong independence between $A,B$ in $\widehat{D}$.} For the sake of contradiction, 
say there is a $(Z,x_i)$-path in $\widehat{D}$. 
Then there exists a $(Z,v)$-path in $D$. This, together with the $(v,X)$-path $P_{\out}$ in $D$, implies an $(Z,X)$-path in $D$, which is a contradiction to the strong independence between the edges of $A, B$ in $D$.

For the sake of contradiction, 
say there is a $(x_i,Z)$-path in $\widehat{D}$. 
Then there exists a $(v,Z)$-path in $D$. This, together with the $(X,v)$-path $P_{\inp}$, implies an $(X,Z)$-path in $D$, which is a contradiction to the strong independence between the edges of $A, B$ in $D$.\\

\noindent{\bf Orderedness.} 
Suppose, for the sake of contradiction, that there is a $(w_j,x_i)$-path in $\widehat{D}$ for $j>i$. Then, there is a $(w_j,v)$-path in $D$. Since $j >i$, this contradicts the choice of $i$.

For the sake of contradiction, say there is a $(x_j,x_i)$-path in $\widehat{D}$ for $j >i$. Then, there is a $(x_j,v)$-path in $D$. Since $j >i$, this contradicts the choice of $i$.

For the sake of contradiction, say there is a $(x_i,x_j)$-path in $\widehat{D}$ for $j <i$. Then, there is a $(v,x_j)$-path in $D$. This together with either the $(x_i,v)$-path or the $(w_i,v)$-path, implies either a $(x_i,x_j)$-path or a $(w_i,x_j)$-path in $D$, contradicting the ordered-ness condition on $A$.\\


\noindent{\bf No bad vertices.} For the sake of contradiction, say there exists a bad vertex $u$ with respect to $(A,B)$ in $\widehat{D}$, that is, $N_{\widehat{D}^{\star}}^-(u) \cap W, N_{\widehat{D}^{\star}}^+(u) \cap X, N_{\widehat{D}^{\star}}^-(u) \cap Y, N_{\widehat{D}^{\star}}^+(u) \cap Z \neq \emptyset$. 
Then, considering that $\widehat{D}$ is obtained from $D$ by only identifying $v$ onto $x_i$, we get that
$N_{D^{\star}}^-(u) \cap W, N_{D^{\star}}^-(u) \cap Y, N_{D^{\star}}^+(u) \cap Z \neq \emptyset$ and $(u,v) \in E(D^{\star})$.  The $(u,v)$-path, together with the $(v,X)$-path $P_{\out}$, implies a $(u,X)$-path in $D$, that is $N_{D^{\star}}^+(u) \cap X \neq \emptyset$. That is, $u$ is a bad vertex with respect to $(A,B)$ in $D$, which is a contradiction.
\end{proof}

Identification Rule~\ref{ir:D} is the analogue of Identification Rule~\ref{ir:C}.

\begin{identification rule}\label{ir:D}
If there exists $v \in V(D) \setminus (W \cup X \cup Y \cup Z)$ such that $N_{D^{\star}}^-(v) \cap Z \neq \emptyset$ and  $N_{D^{\star}}^+(v) \cap Z \neq \emptyset$,
then let $i \in [t]$ be the largest index such that either $(y_i,v) \in E(D^{\star})$ or $(z_i,v) \in E(D^{\star})$. 
Identify $v$ onto $z_i$.
\end{identification rule}

Symmetrically to \Cref{lem:irC-safe}, the following holds.

\begin{lemma}\label{lem:irC-safe}
Let $\widehat{D}$ be the graph obtained after the application of Identification Rule~\ref{ir:C}. Then, the invariants are satisfied.
\end{lemma}

The remaining identification rules are applied on a vertex $v$ when the (in and out) neighbourhood of $v$ is empty in one of the sets $W,X,Y$ or $Z$.

\begin{identification rule}\label{ir:emptyA}
If $N_{D^{\star}}(v) \cap W=\emptyset$ and $\emptyset \neq N_{D^{\star}}^-(v) \cap (X \cup W \cup Y \cup Z) \subseteq Y$, then 
let $i \in [t]$ be the largest integer such that $(y_{i},v) \in E(D^*)$.
Identify $v$ onto $y_{i}$.
\end{identification rule}

\begin{lemma}\label{lem:iremptyA-safe}
Let $\widehat{D}$ be the graph obtained after the application of Identification~\ref{ir:emptyA}. Then, $\widehat{D}$ satisfies the invariants.
\end{lemma}

\begin{proof}
\noindent{\bf Weak independence of $B$ in $\widehat{D}$.} For the sake of contradiction, say there is a $(Z,y_i)$-path in $\widehat{D}$. 
Then there exists a $(Z,v)$-path, in $D$, which is a contradiction (to the assumptions of Identification Rule~\ref{ir:emptyA}).\\

\noindent{\bf Weak independence between $A,B$ in $\widehat{D}$.} For the sake of contradiction, say there is a $(X,y_i)$-path in $\widehat{D}$. 
Then there exists a $(X,v)$-path in $D$, which is a contradiction (to the assumptions of Identification Rule~\ref{ir:emptyA}).\\

\noindent{\bf Strong independence between $A,B$ in $\widehat{D}$.} For the sake of contradiction, 
say there is a $(W,y_i)$-path in $\widehat{D}$. 
Then there exists a $(W,v)$-path in $D$, which is a contradiction (to the assumptions of Identification Rule~\ref{ir:emptyA}). Similarly,  if there is a $(y_i,W)$-path in $\widehat{D}$, then there exists a $(v,W)$-path in $D$, which is again a contradiction.\\

\noindent{\bf Orderedness.} 
Fix $j <i$.
For the sake of contradiction,say there is a $(y_i,y_j)$-path in $\widehat{D}$. Then, there is a $(v,y_j)$-path in $D$. This together with the $(y_i,v)$-path implies a $(y_i,y_j)$-path in $D$, which contradicts the ordered-ness condition on $B$.

Similarly, if there is a $(y_i,z_j)$-path in $\widehat{D}$, then there is a $(v,z_j)$-path in $D$. This, together with the $(y_i,v)$-path, implies a $(y_i,z_j)$-path in $D$, again contradicting the ordered-ness condition on $B$.\\

\noindent{\bf No bad vertices.} For the sake of contradiction, say there exists a bad vertex $u$ with respect to $(A,B)$ in $\widehat{D}$, that is, $N_{\widehat{D}^{\star}}^-(u) \cap W, N_{\widehat{D}^{\star}}^+(u) \cap X, N_{\widehat{D}^{\star}}^-(u) \cap Y, N_{\widehat{D}^{\star}}^+(u) \cap Z \neq \emptyset$. 
Then, considering that $\widehat{D}$ is obtained from $D$ by only identifying $v$ onto $y_i$, we get that
$N_{D^{\star}}^-(u) \cap W, N_{D^{\star}}^+(u) \cap X, N_{D^{\star}}^+(u) \cap Z \neq \emptyset$ and $(v,u) \in E(D^{\star})$.  The $(v,u)$-path, together with the $(v,y_i)$-path, implies a $(u,Y)$-path in $D$, that is $N_{D^{\star}}^-(u) \cap Y \neq \emptyset$. That is, $u$ is a bad vertex with respect to $(A,B)$ in $D$, which is a contradiction.
\end{proof}

The following identification rule is symmetric to \Cref{ir:emptyA}.

\begin{identification rule}\label{ir:emptyB}
If $N_{D^{\star}}(v) \cap Y=\emptyset$ and $\emptyset \neq N_{D^{\star}}^-(v) \cap (X \cup W \cup Y \cup Z) \subseteq W$, then 
let $i \in [t]$ be the largest integer such that $(x_{i},v) \in E(D^*)$.
Identify $v$ onto $w_{i}$.
\end{identification rule}

Similarly to \Cref{lem:iremptyA-safe}, we can prove the following.

\begin{lemma}\label{lem:iremptyB-safe}
Let $\widehat{D}$ be the graph obtained after the application of Identification~\ref{ir:emptyB}. Then, $\widehat{D}$ satisfies the invariants.
\end{lemma}

\begin{identification rule}\label{ir:emptyC}
If $N_{D^{\star}}(v) \cap X=\emptyset$ and $\emptyset \neq N_{D^{\star}}^+(v) \cap (X \cup W \cup Y \cup Z) \subseteq Z$, then 
let $i \in [t]$ be the smallest integer such that $(v,z_i) \in E(D^*)$.
Identify $v$ onto $z_{i}$.
\end{identification rule}

\begin{lemma}\label{lem:iremptyC-safe}
Let $\widehat{D}$ be the graph obtained after the application of Identification~\ref{ir:emptyC}. Then, $\widehat{D}$ satisfies the invariants.
\end{lemma}

\begin{proof}
\noindent{\bf Weak independence of $B$ in $\widehat{D}$.} For the sake of contradiction, say there is a $(z_i,Y)$-path in $\widehat{D}$. 
Then there exists a $(v,Y)$-path, in $D$, which is a contradiction (to the assumptions of Identification Rule~\ref{ir:emptyC}).\\

\noindent{\bf Weak independence between $A,B$ in $\widehat{D}$.} For the sake of contradiction, say there is a $(z_i,W)$-path in $\widehat{D}$. 
Then there exists a $(v,W)$-path in $D$, which is a contradiction (to the assumptions of Identification Rule~\ref{ir:emptyC}).\\

\noindent{\bf Strong independence between $A,B$ in $\widehat{D}$.} For the sake of contradiction, 
say there is a $(z_i,X)$-path in $\widehat{D}$. 
Then there exists a $(v,X)$-path in $D$, which is a contradiction (to the assumptions of Identification Rule~\ref{ir:emptyC}). Similarly,  if there is a $(X,z_i)$-path in $\widehat{D}$, then there exists a $(X,v)$-path in $D$, which is again a contradiction.\\

\noindent{\bf Orderedness.} 
For the sake of contradiction,say there is a $(z_i,z_j)$-path in $\widehat{D}$ for $j <i$.. Then, there is a $(v,z_j)$-path in $D$. This together with the $(z_i,v)$-path implies a $(z_i,z_j)$-path in $D$, which contradicts the ordered-ness condition on $B$.

If there is a $(z_i,y_j)$-path in $\widehat{D}$,
then there is a $(v,y_j)$-path in $D$. This, together with the $(y_i,v)$-path, implies a $(y_i,y_j)$-path in $D$, again contradicting the ordered-ness condition on $B$.\\

\noindent{\bf No bad vertices.} For the sake of contradiction, say there exists a bad vertex $u$ with respect to $(A,B)$ in $\widehat{D}$, that is, $N_{\widehat{D}^{\star}}^-(u) \cap W, N_{\widehat{D}^{\star}}^+(u) \cap X, N_{\widehat{D}^{\star}}^-(u) \cap Y, N_{\widehat{D}^{\star}}^+(u) \cap Z \neq \emptyset$. 
Then, considering that $\widehat{D}$ is obtained from $D$ by only identifying $v$ onto $z_i$, we get that
$N_{D^{\star}}^-(u) \cap W, N_{D^{\star}}^+(u) \cap X, N_{D^{\star}}^-(u) \cap Y \neq \emptyset$ and $(u,v) \in E(D^{\star})$.  The $(u,v)$-path, together with the $(v,z_i)$-path, implies a $(u,Z)$-path in $D$, that is $N_{D^{\star}}^+(u) \cap Z \neq \emptyset$. That is, $u$ is a bad vertex with respect to $(A,B)$ in $D$, which is a contradiction.
\end{proof}

The following identification rule is symmetric to \Cref{ir:emptyC}.

\begin{identification rule}\label{ir:emptyD}
If $N_{D^{\star}}(v) \cap Z=\emptyset$ and $\emptyset \neq N_{D^{\star}}^+(v) \cap (X \cup W \cup Y \cup Z) \subseteq X$, then 
let $i \in [t]$ be the smallest integer such that $(v,x_i) \in E(D^*)$.
Identify $v$ onto $x_{i}$.
\end{identification rule}

Similarly to \Cref{lem:iremptyC-safe}, we can prove the following.

\begin{lemma}\label{lem:iremptyD-safe}
Let $\widehat{D}$ be the graph obtained after the application of Identification~\ref{ir:emptyD}. Then, $\widehat{D}$ satisfies the invariants.
\end{lemma}

\begin{lemma}\label{lem:ordered-cleaning-exhaustive}
When none of the Identification Rules~\ref{ir:A}-\ref{ir:emptyD} are applicable, for each vertex of $v \in V(D) \setminus V(A\cup B)$ either $N_D^-(v) =\emptyset$ or $N_D^+(v) = \emptyset$.
\end{lemma}
\begin{proof}
To prove the lemma, we prove the following claims.

\begin{claim}\label{claim:ordered-ir}
If $N_{D^{\star}}^+(v) \cap W \neq \emptyset$, then
$ N_{D^{\star}}^-(v) \cap (W \cup X \cup Y \cup Z) \subseteq W$. 
\end{claim}

\begin{claimproof}
If $x_j \in N_{D^{\star}}^-(v) $ for some $j \in [t]$, then this implies a $(X,W)$-path in $D$, contradicting the weak independence of $A$.
If $y_j \in N_{D^{\star}}^-(v)$ for some $j \in [t]$, then this implies a $(Y,W)$-path in $D$, contradicting the strong independence between $A,B$. 
If $z_j \in N_{D^{\star}}^-(v)$ for some $j \in [t]$, then this implies a $(Z,W)$-path in $D$, contradicting the weak independence between $A,B$. 
\end{claimproof}

Symmetrically to \Cref{claim:ordered-ir}, the following holds.

\begin{claim}\label{claim:ordered-ir2}
If $N_{D^{\star}}^+(v) \cap Y \neq \emptyset$, then
$ N_{D^{\star}}^-(v) \cap (W \cup X \cup Y \cup Z) \subseteq Y$. 
\end{claim}

\begin{claim}\label{claim:ordered-ir3}
If $N_{D^{\star}}^-(v) \cap X \neq \emptyset$, then
$ N_{D^{\star}}^+(v) \cap (W \cup X \cup Y \cup Z) \subseteq X$. 
\end{claim}

\begin{claimproof}
If $w_j \in N_{D^{\star}}^+(v) $ for some $j \in [t]$, then this implies a $(X,W)$-path in $D$, contradicting the weak independence of $A$.
If $y_j \in N_{D^{\star}}^+(v)$ for some $j \in [t]$, then this implies a $(X,Y)$-path in $D$, contradicting the weak independence between $A,B$. 
If $z_j \in N_{D^{\star}}^+(v)$ for some $j \in [t]$, then this implies a $(X,Z)$-path in $D$, contradicting the strong independence between $A,B$. 
\end{claimproof}

Symmetrically to \Cref{claim:ordered-ir3}, the following holds.

\begin{claim}\label{claim:ordered-ir4}
If $N_{D^{\star}}^-(v) \cap Z \neq \emptyset$, then
$ N_{D^{\star}}^+(v) \cap (W \cup X \cup Y \cup Z) \subseteq Z$. 
\end{claim}

From Claims~\ref{claim:ordered-ir} and~\ref{claim:ordered-ir4}, if Identification Rules~\ref{ir:A}-\ref{ir:D} are not applicable and there exists $v \in W \cup X \cup Y \cup Z$ such that $N_{D^{\star}}^+(v) \cap (W \cup X \cup Y \cup Z) \neq \emptyset$ and $N_{D^{\star}}^-(v) \cap (W \cup X \cup Y \cup Z) \neq \emptyset$, then since there are no bad vertices with respect to $(A,B)$, we conclude that either $N_{D^{\star}}(v) \cap W =\emptyset$ or $N_{D^{\star}}(v) \cap X =\emptyset$ or $N_{D^{\star}}(v) \cap Y=\emptyset$ or $N_{D^{\star}}(v) \cap Z =\emptyset$. Thus, Claims~\ref{claim:ordered-ir5}-\ref{claim:ordered-ir8} prove the lemma.

\begin{claim}\label{claim:ordered-ir5}
If Identification Rules~\ref{ir:A}-\ref{ir:D} are no longer applicable, $N_{D^{\star}}(v) \cap W=\emptyset$ and  $N_{D^{\star}}^+(v) \cap (X \cup Y \cup Z) \neq \emptyset$, then $N_{D^{\star}}^-(v) \cap (X \cup Y \cup Z) \subseteq Y$.
\end{claim}

\begin{claimproof}
We first show that that $N_{D^{\star}}^-(v) \cap X = \emptyset$. 
For the sake of contradiction, say $N_{D^{\star}}^-(v) \cap X \neq \emptyset$. 
Then, since Identification Rule~\ref{ir:B} is not applicable, $N_{D^{\star}}^+(v) \cap X = \emptyset$.  Also, $N_{D^{\star}}^+(v) \cap Y = \emptyset$, as otherwise there is a $(X,Y)$-path in $D$, contradicting the weak independence between $A,B$. Further, $N_{D^{\star}}^+(v) \cap Z = \emptyset$, as otherwise there is a $(X,Z)$-path in $D$, contradicting the strong independence between $A,B$.

Similarly, if $N_{D^{\star}}^-(v) \cap Z \neq \emptyset$. 
Then, since Identification Rule~\ref{ir:B} is not applicable, $N_{D^{\star}}^+(v) \cap Z = \emptyset$.  Also, 
$N_{D^{\star}}^+(v) \cap X = \emptyset$, as otherwise there is a $(Z,X)$-path in $D$, contradicting the strong independence between $A,B$. Further, 
$N_{D^{\star}}^+(v) \cap Y = \emptyset$, as otherwise there is a $(Z,Y)$-path in $D$, contradicting the weak independence of $B$.
\end{claimproof}

Similarly to \Cref{claim:ordered-ir5}, we can prove the following.

\begin{claim}\label{claim:ordered-ir6}
If Identification Rules~\ref{ir:A}-\ref{ir:D} are no longer applicable, $N_{D^{\star}}(v) \cap Y=\emptyset$ and  $N_{D^{\star}}^+(v) \cap (X \cup W \cup Z) \neq \emptyset$, then $N_{D^{\star}}^-(v) \cap (W \cup X \cup Z) \subseteq W$.
\end{claim}

\begin{claim}\label{claim:ordered-ir7}
If Identification Rules~\ref{ir:A}-\ref{ir:D} are no longer applicable, $N_{D^{\star}}(v) \cap X=\emptyset$ and $N_{D^{\star}}^-(v) \cap (X \cup Y \cup Z) \neq \emptyset$, then $N_{D^{\star}}^+(v) \cap ( W \cup Y \cup Z) \subseteq Z$.
\end{claim}

\begin{claimproof}
We first show that that $N_{D^{\star}}^+(v) \cap W = \emptyset$. 
For the sake of contradiction, say $N_{D^{\star}}^+(v) \cap W \neq \emptyset$. 
Then, since Identification Rule~\ref{ir:B} is not applicable, $N_{D^{\star}}^-(v) \cap W = \emptyset$. 
Also, $N_{D^{\star}}^-(v) \cap Y = \emptyset$, as otherwise there is a $(Y,W)$-path in $D$, contradicting the strong independence between $A,B$. 
Further, $N_{D^{\star}}^-(v) \cap Z = \emptyset$, as otherwise there is a $(Z,W)$-path in $D$, contradicting the weak independence between $A,B$.

Similarly, if $N_{D^{\star}}^+(v) \cap Y \neq \emptyset$,
then since Identification Rule~\ref{ir:B} is not applicable, $N_{D^{\star}}^-(v) \cap Y = \emptyset$.  
Also,  $N_{D^{\star}}^-(v) \cap W= \emptyset$, as otherwise there is a $(W,Z)$-path in $D$, contradicting the strong independence between $A,B$. 
Further, 
$N_{D^{\star}}^-(v) \cap Z= \emptyset$, as otherwise there is a $(Z,Y)$-path in $D$, contradicting the weak independence of $B$.
\end{claimproof}

Similarly to \Cref{claim:ordered-ir7}, we can prove the following.

\begin{claim}\label{claim:ordered-ir8}
If Identification Rules~\ref{ir:A}-\ref{ir:D} are no longer applicable, $N_{D^{\star}}(v) \cap Z=\emptyset$ and $N_{D^{\star}}^-(v) \cap (X \cup Y \cup Z) \neq \emptyset$, then $N_{D^{\star}}^+(v) \cap ( W \cup X \cup Y) \subseteq X$.
\end{claim}
This concludes the proof of \Cref{lem:ordered-cleaning-exhaustive}.
\end{proof}

Now suppose that Identification Rules~\ref{ir:A}-\ref{ir:emptyD} are no longer applicable. From Lemma~\ref{lem:ordered-cleaning-exhaustive}, each vertex of $V(D) \setminus (W \cup X \cup Y \cup Z)$ either has only in-neighbours or only out-neighbours in $D^{\star}$.
Let $V_{\sss} =\{v : v \in V(D) \setminus (W \cup X \cup Y \cup Z) \text{ and } N_{D^{\star}}^+(v) \cap (W \cup X \cup Y \cup Z) \neq \emptyset \}$ and
$V_{\ttt} = V(D) \setminus (W \cup X \cup Y \cup Z \cup V_{\sss})$. Then identify all the vertices of $V_{\sss}$ onto a new vertex $\sss$ and all the vertices of $V_{\ttt}$ onto a new vertex $\ttt$. The resulting graph is a semi-cleaned ordered $t$-\toughpair.

\subsubsection{From semi-cleaned ordered \toughpair\ to hard-pattern}
\label{sec:semicleantoclean}

The aim of this section is to prove \Cref{lem:semicleantoclean}. We first introduce the \emph{almost $t$-hard matching patterns} which differ from the $t$-hard matching patterns in that the source and sink vertices of items 3 and 4, respectively, in \Cref{def:cleanedorderedtoughpair} may not have "full" neighborhoods (recall that an $S$-source is a source vertex~$s$ such that $N^+(s) = S$ and that an $S$-sink is a sink vertex such that $N^-(s) = S$).

\begin{definition}[almost $t$-hard matching pattern]
\label{def:almost}
An \emph{almost $t$-hard matching pattern} is an (acyclic) digraph $D$ constructed the following way. We start with disjoint vertex sets $W = \{w_1,\ldots,w_t\}$, $X = \{x_1,\ldots,x_t\}$, $Y = \{y_1,\ldots,y_t\}$ and $Z = \{z_1,\ldots,z_t\}$ and introduce the edges $(w_i,x_i)$ and $(y_i,z_i)$ for every $i\in[t]$.
Furthermore, we introduce into $D$ any combination of the following items:
\begin{enumerate}
\item[1.] either the directed path $w_1 \rightarrow w_2 \rightarrow \ldots \rightarrow w_t \rightarrow z_1 \rightarrow z_2 \rightarrow \ldots \rightarrow z_t$, or any of the directed paths $w_1 \rightarrow w_2 \rightarrow \ldots \rightarrow w_t$ and 
$z_1 \rightarrow z_2 \rightarrow \ldots \rightarrow z_t$;
\item[2.] either the directed path $y_1 \rightarrow y_2 \rightarrow \ldots \rightarrow y_t \rightarrow x_1 \rightarrow x_2 \rightarrow \ldots \rightarrow x_t$, or any of the directed paths $x_1 \rightarrow x_2 \rightarrow \ldots \rightarrow x_t$ and $y_1 \rightarrow y_2 \rightarrow \ldots \rightarrow y_t$;
\item[3.] a vertex~$s$ such that $N^-(s) = \emptyset$ and $N^+(s) \subseteq W \cup X \cup Y \cup Z$; 
\item[4.] a vertex $t$ such that $N^+(t) = \emptyset$ and $N^-(t) \subseteq W \cup X \cup Y \cup Z$; 
\item[5.] a vertex $r_{WZ}$ such that $N^-(r_{WZ}) = W$ and $N^+(r_{WZ}) = Z$;
\item[6.] a vertex $r_{YX}$ such that $N^-(r_{YX}) = Y$ and $N^+(r_{YX}) = X$.
\end{enumerate}
\end{definition}

We further introduce the four canonical graphs below corresponding to "half" a hard matching pattern (see \Cref{fig:conf} for an illustration).

\begin{definition}
\label{def:canonical}
Let $U = \{u_1,\ldots,u_t\}$ and $V = \{v_1,\ldots,v_t\}$. We define the following four reachability-minimal acyclic digraphs on vertex set $U \cup V$.
\begin{itemize}
\item The digraph $O_1^t$ is a $(U,V)$-induced-matching.
\item The digraph $O_2^t$ is a $(U,V)$-matching which further contains the directed path on $U$ from $u_1$ to $u_t$.
\item The digraph $O_3^t$ is a $(U,V)$-matching which further contains the directed path on $V$ from $v_1$ to $v_t$.
\item The digraph $O_4^t$  is a $(U,V)$-matching which further contains the directed path on $U$ from $u_1$ to $u_t$ and the directed path on $V$ from $v_1$ to $v_t$.
\end{itemize}
\end{definition}

\newcommand{\onepath}[4]{
\foreach \i in {1,...,4}
{
\pgfmathsetmacro{\x}{0.75*\i+#1}
\pgfmathsetmacro{\y}{#2+.75}
\node[circ] (a\i) at (\x,\y) {};
\node[circ] (b\i) at (\x,#2) {};
\draw[->,>=stealth,thick] (a\i) -- (b\i);
}

\ifthenelse{#3=1}{
\foreach \i in {1,...,3}
{
\pgfmathtruncatemacro{\j}{\i+1}
\draw[->,>=stealth,thick] (a\i) -- (a\j);
}
}{}

\ifthenelse{#4=1}{
\foreach \i in {1,...,3}
{
\pgfmathtruncatemacro{\j}{\i+1}
\draw[->,>=stealth,thick] (b\i) -- (b\j);
}
}{}
}

\begin{figure}
\centering
\begin{subfigure}[h]{.23\textwidth}
\centering
\begin{tikzpicture}
\onepath{0}{0}{0}{0};
\end{tikzpicture}
\caption*{$O_1^4$}
\end{subfigure}
\begin{subfigure}[h]{.23\textwidth}
\centering
\begin{tikzpicture}
\onepath{0}{0}{1}{0};
\end{tikzpicture}
\caption*{$O_2^4$}
\end{subfigure}
\begin{subfigure}[h]{.23\textwidth}
\centering
\begin{tikzpicture}
\onepath{0}{0}{0}{1};
\end{tikzpicture}
\caption*{$O_3^4$}
\end{subfigure}
\begin{subfigure}[h]{.23\textwidth}
\centering
\begin{tikzpicture}
\onepath{0}{0}{1}{1};
\end{tikzpicture}
\caption*{$O_4^4$}
\end{subfigure}
\caption{The graphs of \Cref{def:canonical} for $t=4$.}
\label{fig:conf}
\end{figure}

\semicleantoclean*

\begin{proof}
For every $p \geq 1$ and $j \geq 1$, let $\ell_j(p)$ be an integer such that $p \geq \mathcal{R}(\ell_j(p),j)$.
For notational convenience, we show, instead of the above statement, that if $D$ is a semi-cleaned ordered $t$-\toughpair\ then either $D^*$ contains a $g(t)$-induced-biclique whose edges are minimal in $D$, or $D$ can be identified to a digraph $\hat{D}$ such that $\hat{D}$ is transitively equivalent to a $g(t)$-hard matching pattern, where $g(t) = (\ell_4(\ell_4(\ell_5(t))/2)-4)/8$ (equivalently, $h(t) = \mathcal{R}(2\mathcal{R}(2\mathcal{R}(8t+4,4),4),5)$).

Let $D$ be a semi-cleaned ordered $t$-\toughpair\ and let $(A,B)$ be the ordered $t$-\toughpair\ contained in $D$,
where $A = ((w_1,x_1),\ldots,(w_t,x_t))$ and $B = ((y_1,z_1),\ldots,(y_t,z_t))$.
Denote by $W= \{w_i~|~i \in [t]\}$, $X = \{x_i~|~i \in [t]\}$, $Y= \{y_i~|~i \in [t]\}$ and $Z= \{z_i~|~i\in [t]\}$.
Given a digraph $H$, we denote by $\overline{H}$ the underlying undirected graph and by $H^*$ the transitive closure of $H$. Before turning to the proof of the lemma, we first show the following claims. 

\begin{claim}
\label{clm:ramseyA}
One of the following holds.
\begin{itemize}
\item[(1)] There exist subsets $X' \subseteq X$ and $W' \subseteq W$ such that $|X'| = |W'| = \lfloor \ell_5(t)/2 \rfloor$ and $D$ contains a $(W',X')$-induced-biclique.
\item[(2)] There exist a set $I \subseteq [t]$ of size $\ell_5(t)$ such that $D[\{x_i,w_i~|~i \in I\}]$ is transitively equivalent to one of $O_1^{\ell_5(t)},O_2^{\ell_5(t)},O_3^{\ell_5(t)}$ and $O_4^{\ell_5(t)}$.
\end{itemize}
\end{claim}

\begin{claimproof}
Consider the edge-coloring of the complete undirected graph $G$ on vertex set $\{1,\ldots,t\}$ defined as follows: for every $i,j \in [t]$ where $i < j$, the edge $ij \in E(G)$ receives 
\begin{itemize}
\item color 1 if $(w_i,w_j),(x_i,x_j),(w_i,x_j) \notin E(D^*)$;
\item color 2 if $(w_i,w_j) \in E(D)$ and $(x_i,x_j) \notin E(D^*)$;
\item color 3 if $(w_i,w_j) \notin E(D)$ and $(x_i,x_j) \in E(D^*)$;
\item color 4 if $(w_i,w_j),(x_i,x_j) \in E(D^*)$; and
\item color 5 if $(w_i,w_j),(x_i,x_j) \notin E(D^*)$ and $(w_i,x_j) \in E(D^*)$.
\end{itemize}
Note that for any $1 \leq i<j \leq t$, if $D^*$ contains one of $(w_i,w_j)$ and $(x_i,x_j)$ then $D^*$ also contains $(w_i,x_j)$ and thus, each edge of $E(G)$ receives a color in the above coloring.
Now by Ramsey's Theorem, $G$ contains a monochromatic clique of size $\ell_5(t)$: let $1 \leq i_1 < \ldots < i_{\ell_5(t)} \leq t$ be a set of $\ell_5(t)$ vertices inducing a monochromatic clique in $G$. Since $(A,B)$ is an ordered $t$-\toughpair, for every $1 \leq p < q \leq \ell_5(t)$, there is no $(w_{i_q},w_{i_p})$-path, no $(w_{i_q},x_{i_p})$-path and no $(x_{i_q},x_{i_p})$-path. It follows that if the clique in $G$ on vertex set $\{i_1,\ldots,i_{\ell_5(t)}\}$ has 
\begin{itemize}
\item color 1 then $D[\{w_{i_j},x_{i_j}~|~ j \in [\ell_5(t)]\}]$ is transitively equivalent to $O^{\ell_5(t)}_1$;
\item color 2 then $D[\{w_{i_j},x_{i_j}~|~ j \in [\ell_5(t)]\}]$ is transitively equivalent to $O^{\ell_5(t)}_2$;
\item color 3 then $D[\{w_{i_j},x_{i_j}~|~ j \in [\ell_5(t)]\}]$ is transitively equivalent to $O^{\ell_5(t)}_3$;
\item color 4 then $D[\{w_{i_j},x_{i_j}~|~ j \in [\ell_5(t)]\}]$ is transitively equivalent to $O^{\ell_5(t)}_4$;
\item color 5 then $D$ contains a $(\{w_{i_j}~|~j \in [\lfloor \ell_5(t)/2 \rfloor]\},\{x_{i_{\lfloor \ell_5(t)/2 \rfloor + j}}~|~j \in [\lfloor \ell_5(t)/2 \rfloor]\})$-induced-biclique.
\end{itemize}
Thus, if the clique in $G$ on vertex set $\{i_1,\ldots,i_{\ell_5(t)}\}$ does not have color 5 then we may take $I = \{i_1,\ldots,i_{\ell_5(t)}\}$ to prove our claim.  
\end{claimproof}

Symmetrically to \Cref{clm:ramseyA}, we have the following.

\begin{claim}
\label{clm:ramseyB}
One of the following holds.
\begin{itemize}
\item[(1)] There exist subsets $Y' \subseteq Y$ and $Z' \subseteq Z$ such that $|Y'| = |Z'| = \lfloor \ell_5(t)/2 \rfloor$ and $D$ contains a $(Y',Z')$-induced-biclique.
\item[(2)] There exist a set $I \subseteq [t]$ of size $\ell_5(t)$ such that $D[\{y_i,z_i~|~i \in I\}]$ is transitively equivalent to one of $O_1^{\ell_5(t)},O_2^{\ell_5(t)},O_3^{\ell_5(t)}$ and $O_4^{\ell_5(t)}$.
\end{itemize}
\end{claim}

\begin{claim}
\label{clm:ramseyWZ}
For any subsets $W' \subseteq W$ and $Z' \subseteq Z$ where $|W'| = |Z'| = p$, there exist subsets 
$W'' \subseteq W'$ and $Z'' \subseteq Z'$ such that $|W''| = |Z''| = \lfloor \ell_4(p) / 2 \rfloor$ and
\begin{itemize}
\item[$\bullet$] either there is no edge in $D$ from a vertex of $W''$ to a vertex of $Z''$,
\item[$\bullet$] or $D^*$ contains a $(W'',Z'')$-biclique (not necessarily induced).
\end{itemize}
\end{claim}

\begin{claimproof}
Let $W' \subseteq W$ and $Z' \subseteq Z$ be two subsets of size $p$. We distinguish two cases depending on whether the bipartite graph $\overline{D}[W',Z']$ has a matching of size at least $p/2$ or not.

\paragraph{Case 1. The bipartite graph $\overline{D}[W',Z']$ has a matching of size at least $p/2$.} Let $M$ be a matching in $\overline{D}[W',Z']$ of size $\lfloor p/2 \rfloor$ and let $e_1,\ldots,e_{\lfloor p/2 \rfloor}$ be an arbitrary ordering of the edges in $M$. Consider the edge-coloring of the complete undirected graph $G$ on vertex set $\{1,\ldots, \lfloor p/2 \rfloor\}$ defined as follows: for every $1 \leq i < j \leq \lfloor p/2 \rfloor$, the edge $ij \in E(G)$ receives
\begin{itemize}
\item color 1 if $(\tail(e_i),\head(e_j)),(\tail(e_j),\head(e_i)) \notin E(D^*)$;
\item color 2 if $(\tail(e_i),\head(e_j)) \in E(D^*)$ and $(\tail(e_j),\head(e_i)) \notin E(D^*)$;
\item color 3 if $(\tail(e_i),\head(e_j)) \notin E(D^*)$ and $(\tail(e_j),\head(e_i)) \in E(D^*)$; and
\item color 4 if $(\tail(e_i),\head(e_j)),(\tail(e_j),\head(e_i)) \in E(D^*)$;
\end{itemize}
Then by Ramsey's Theorem, $G$ contains a monochromatic clique of size $\ell_4(p)$: let $1 \leq i_1 < \ldots < i_{\ell_4(p)} \leq p/2$ be a set of $\ell_4(p)$ vertices inducing a monochromatic clique in $G$. Now if this clique has
\begin{itemize}
\item color 1 then by taking $W'' = \{\tail(e_{i_j})~|~j \in [\lfloor \ell_4(p)/2 \rfloor]\}$ and $Z'' = \{\head(e_{i_j})~|~j \in [\lfloor \ell_4(p)/2 \rfloor]\}$, there is no edge from $W''$ to $Y''$ in $D^*$ (and a fortiori in $D$);
\item color 2 then by taking $W'' = \{\tail(e_{i_j})~|~j \in [\lfloor \ell_4(p)/2 \rfloor]\}$ and $Z'' = \{\head(e_{i_{\lfloor \ell_4(p)/2 \rfloor + j}})~|~j \in [\lfloor \ell_4(p)/2 \rfloor]\}$, $D^*$ contains a $(W'',Z'')$-biclique;
\item color 3 then by taking $W'' = \{\tail(e_{i_{\lfloor \ell_4(p)/2 \rfloor + j}})~|~j \in [\lfloor \ell_4(p)/2 \rfloor]\}$ and $Z'' = \{\head(e_{i_j})~|~j \in [\lfloor \ell_4(p)/2 \rfloor]\}$, $D^*$ contains a $(W'',Z'')$-biclique;
\item color 4 then by taking $W'' = \{\tail(e_{i_j})~|~j \in [\lfloor \ell_4(p)/2 \rfloor]\}$ and $Z'' = \{\head(e_{i_j})~|~j \in [\lfloor \ell_4(p)/2 \rfloor]\}$, $D^*$ contains a $(W'',Z'')$-biclique.
\end{itemize}

\paragraph{Case 2. The bipartite graph $\overline{D}[W',Z']$ has no matching of size at least $p/2$.} Then by K\"onig's Theorem, $\overline{D}[W',Z']$ has, in this case, a vertex cover at size at most $p/2$: let $V \subseteq W' \cup Z'$ be a minimum vertex cover of $\overline{D}[W',Z']$. Since then, $\min \{ W' \setminus V,Z' \setminus V\} \geq \lfloor \ell_4(p)/ 2 \rfloor \geq \lfloor \ell_4(p)/2 \rfloor$, we may take any subsets $W'' \subseteq W' \setminus V$ and $Z'' \subseteq Z' \setminus V$ of size $\lfloor \ell_4(p)/2 \rfloor$ to prove the claim, as there are surely no edges from $W''$ to $Z''$ in $D$. 
\end{claimproof}

Symmetrically to \Cref{clm:ramseyWZ}, we have the following.

\begin{claim}
\label{clm:ramseyXY}
For any subsets $Y' \subseteq Y$ and $X' \subseteq X$ where $|Y'| = |X'| = p$, there exist subsets 
$Y'' \subseteq Y'$ and $X'' \subseteq X'$ such that $|Y''| = |X''| = \lfloor \ell_4(p)/2 \rfloor$ and
\begin{itemize}
\item[$\bullet$] either there is no edge in $D$ from a vertex of $Y''$ to a vertex of $X''$,
\item[$\bullet$] or $D^*$ contains a $(Y'',X'')$-biclique (not necessarily induced).
\end{itemize}
\end{claim}

We next define partitions with respect to a set and prove thereafter useful properties of these partitions.

\begin{definition}
For any $V \in \{W,X,Y,Z\}$ and any $U \subseteq V$, the \emph{partition of $V$ w.r.t. $U$} is the partition of $V \setminus U$ into four sets $V^\star, V^\circ, V^+$ and $V^-$ such that the following hold.
\begin{itemize}
\item For every $v \in V^\star$, $N^+_{D^*}(v) \cap U \neq \emptyset$ and $N^-_{D^*}(v) \cap U \neq \emptyset$.
\item For every $v \in V^\circ$, $(N^+_{D^*}(v) \cup N^-_{D^*}(v)) \cap U = \emptyset$.
\item For every $v \in V^+$, $N^+_{D^*}(v) \cap U \neq \emptyset$ and $N^-_{D^*}(v) \cap U = \emptyset$.
\item For every $v \in V^-$, $N^+_{D^*}(v) \cap U  = \emptyset$ and $N^-_{D^*}(v) \cap U \neq \emptyset$.
\end{itemize} 
\end{definition}

Note that since $(A,B)$ is an ordered \toughpair, for any $V \in \{W,X,Y,Z\}$ and any $v \in V$, 
\begin{itemize}
\item if $u \in N^-_{D^*}(v)$ then $u$ has a smaller index than that of $v$, and
\item if $u \in N^+_{D^*}(v)$ then $u$ has a larger index than that of $v$. 
\end{itemize}

\begin{claim}
\label{clm:independent}
For any $V \in \{W,X,Y,Z\}$ and any $U \subseteq V$, the following hold.
Let $V^\star, V^\circ, V^+$ and $V^-$ be the partition of $V$ w.r.t $U$. Then there is no edge in $D^*$
\begin{itemize}
\item[(i)] from a vertex of $V^\star$ to a vertex of $V^\circ \cup V^+$, 
\item[(ii)] from a vertex of $V^-$ to a vertex of $V^\star \cup V^\circ \cup V^+$, and 
\item[(iii)] from a vertex of $V^\circ$ to a vertex of $V^+ \cup V^\star$.
\end{itemize}
Furthermore, if $D^*[U]$ has no edge then $V^\star = \emptyset$.
\end{claim}

\begin{claimproof}
Consider $u \in V^\star$. If there exists $v \in V^\circ \cup V^+$ such that $(u,v) \in E(D^*)$ then in particular $\emptyset \neq N^-_{D^*}(u) \cap U \subseteq N^-_{D^*}(v)$, a contradiction to the fact that $v \in V^+ \cup V^\circ$.
Similarly, there is no edge in $D^*$ from a vertex of $V^-$ to a vertex of $V^+ \cup V^\star$: indeed, if there exist $u \in V^-$ and $v\in V^+ \cup V^\star$ such that $(u,v) \in E(D^*)$ then in particular $\emptyset \neq N^-_{D^*}(u) \cap U \subseteq N^-_{D^*}(v)$, a contradiction to the fact that $v \in V^+$. Furthermore, there is no edge from a vertex of $V^-$ to a vertex of $V^\circ$ for if there exist $u \in V^-$ and $v \in V^\circ$ such that $(u,v) \in E(D^*)$, then in particular $\emptyset \neq N^-_{D^*}(u) \cap U \subseteq N^-_{D^*}(v)$, a contradiction to the fact that $v \in V^\circ$. Finally, if there exist $u \in V^\circ$ and $v \in V^+ \cup V^\star$ such that $(u,v) \in E(D^*)$ then in particular $\emptyset \neq N^+_{D^*}(v) \cap U \subseteq N^+_{D^*}(u)$, a contradiction to the fact that $u \in V^\circ$.
Now assume that $D^*[U]$ has no edge and suppose for a contradiction that there exists $u \in V \setminus U$ such that both $N^-_{D^*}(u) \cap U \neq \emptyset$ and $N^+_{D^*}(u) \cap U \neq \emptyset$. Then for any $x \in N^-_{D^*}(u) \cap U$ and $y \in N^+_{D^*}(u) \cap U$, $(x,y) \in E(D^*)$, a contradiction to the fact that $U$ is an independent set of $D^*$. 
\end{claimproof}

\begin{claim}
\label{clm:W+}
For any $V \in \{W,Y\}$ and any $U \subseteq V$, the following hold:
\begin{itemize}
\item if there exist $u \in W \cup X \cup Y \cup Z$ and $v \in V^+$ such that $(u,v) \in E(D^*)$ then $u \in V^+$;
\item if there exist $u \in W \cup X \cup Y \cup Z$ and $v \in V^\circ$ such that $(u,v) \in E(D^*)$ then $u \in V^\circ \cup V^+$,
\end{itemize}
where $V^\star, V^\circ, V^+,V^-$ is the partition of $V$ w.r.t $U$.
\end{claim}

\begin{claimproof}
Suppose that there exist $u \in W \cup X \cup Y \cup Z$ and $v \in V^+$ such that $(u,v) \in E(D^*)$. Observe first that since $V \in \{W,Y\}$ and $(A,B)$ is an ordered \toughpair, necessarily $u \in V$. Now since $v \in V^+$, $u \notin U$ by definition and since by \Cref{clm:independent}, $u \notin V^\star \cup V^\circ \cup V^-$, we conclude that $u \in V^+$.
Suppose next that there exist $u \in W \cup X \cup Y \cup Z$ and $v \in V^\circ$ such that $(u,v) \in E(D^*)$. Then as previously $u \in V$. Now since $v \in V^\circ$, $u \notin U$ by definition and since by \Cref{clm:independent}, $u \notin V^\star \cup V^-$, we conclude that $u \in V^\circ \cup V^+$.
\end{claimproof}

\begin{claim}
\label{clm:X-}
For any $V \in \{X,Z\}$ and any $U \subseteq V$, the following hold:
\begin{itemize}
\item if there exist $u \in V^-$ and $v \in W \cup X \cup Y \cup Z$ such that $(u,v) \in E(D^*)$ then $v \in V^-$;
\item if there exist $v \in V^\circ$ and $v \in W \cup X \cup Y \cup Z$ such that $(u,v) \in E(D^*)$ then $u \in V^\circ \cup V^-$,
\end{itemize}
where $V^\star, V^\circ, V^+,V^-$ is the partition of $V$ w.r.t $U$. 
\end{claim}

\begin{claimproof}
Suppose that there exist $u \in V^-$ and $v \in W \cup X \cup Y \cup Z$ such that $(u,v) \in E(D^*)$. Observe first that since $V \in\{X,Z\}$ and $(A,B)$ is an ordered \toughpair, necessarily $v \in V$. Now since $u \in V^-$, $v \notin U$ by definition and since by \Cref{clm:independent}, $v \notin V^\star \cup V^\circ \cup V^+$, we conclude that $v \in V^-$.
Suppose next that there exist $v \in V^\circ$ and $v \in W \cup X \cup Y \cup Z$ such that $(u,v) \in E(D^*)$ then $u \in V^\circ \cup V^-$. Then as previously $v \in V$. Now since $u \in V^\circ$,  $v \notin U$ by definition and since by \Cref{clm:independent}, $v \notin V^+ \cup V^\star$, we conclude that $v \in V^\circ \cup V^-$.
\end{claimproof}

\newcommand{\sizeWZ}{\ell_4(\ell_5(t))/2}
\newcommand{\sizeXY}{4g(t)+2}
\newcommand{\sizeXYo}{4g(t)+1}
\newcommand{\sizeXYt}{4g(t)}

We now turn to the proof of the lemma. We assume henceforth that $D$ contains no $g(t)$-induced-biclique (we are done otherwise). We first prove that $D$ can be identified to a digraph which is transitively equivalent to an almost $\sizeXYt$-hard matching pattern and then show how to obtain a $g(t)$-hard matching pattern from an almost $4g(t)$-hard matching pattern. 

By \Cref{clm:ramseyA}, there exists a set $I_A \subseteq [t]$ of size $\ell_5(t)$ such that $D[\{w_i,x_i~|~i \in I_A\}]$ is transitively equivalent to one of $O_1^{\ell_5(t)},O_2^{\ell_5(t)},O_3^{\ell_5(t)}$ and $O_4^{\ell_5(t)}$; and similarly, by \Cref{clm:ramseyB}, there exists a set $I_B \subseteq [t]$ of size $\ell_5(t)$ such that $D[\{y_i,z_i~|~i\in I_B\}]$ is transitively equivalent to one of $O_1^{\ell_5(t)},O_2^{\ell_5(t)},O_3^{\ell_5(t)}$ and $O_4^{\ell_5(t)}$. 
Now by \Cref{clm:ramseyWZ} there exist subsets $I_A^1 \subseteq I_A$ and $I_B^1 \subseteq I_B$ such that $|I_A^1| = |I_B^1| = \sizeWZ$ and there is either no edge or every edge in $D$ from $\{w_i~|~i \in I_A^1\}$ to $\{z_i~|~i \in I_B^1\}$. 
By \Cref{clm:ramseyXY}, there further exist subsets $I_A^2 \subseteq I_A^1$ and $I_B^2 \subseteq I_B^1$ such that $|I_A^2| = |I_B^2| = \ell_4(\ell_4(\ell_5(t))/2))/2 = \sizeXY$ and there is either no edge or every edge in $D$ from $\{y_i~|~i \in I_A^2\}$ to $\{x_i~|~i \in I_B^2\}$.
 
Now let $W^\star,W^\circ,W^+,W^-$ be the partition of $W$ w.r.t. $\{w_i~|~i \in I_A^2\}$,
$X^\star,X^\circ,X^+,X^-$ be the partition of $X$ w.r.t $\{x_i~|~i \in I_A^2\}$,
$Y^\star,Y^\circ,Y^+,Y^-$ be the partition of $Y$ w.r.t $\{y_i~|~i \in I_B^2\}$ 
and $Z^\star,Z^\circ,Z^+,Z^-$ be the partition of $Z$ w.r.t $\{z_i~|~i \in I_B^2\}$.
We devise the following identification rules.

\begin{identification rule}
\label{ir:source}
Proceed as follows.
\begin{enumerate}[nosep]
\item Identify every vertex in $W^\circ \cup W^+ \cup Y^\circ \cup Y^+$ (together with the source vertex $\sss$ of $D$ when it exists) to a single vertex $s$.
\item If $D[\{w_i,x_i~|~i \in I_A\}]$ is transitively equivalent to $O_1^{\ell_5(t)}$, further identity every vertex in $X^+$ to $s$.
\item If $D[\{y_i,z_i~|~i \in I_B\}]$ is transitively equivalent to $O_1^{\ell_5(t)}$, further identity every vertex in $Z^+$ to $s$.
\end{enumerate}
\end{identification rule}

\begin{identification rule}
\label{ir:sink}
Proceed as follows.
\begin{enumerate}[nosep]
\item Identify every vertex in $X^\circ \cup X^- \cup Z^\circ \cup Z^-$ (together with the sink vertex $\ttt$ of $D$ when it exists) to a single vertex $t$.
\item If $D[\{w_i,x_i~|~i \in I_A\}]$ is transitively equivalent to $O_1^{\ell_5(t)}$, further identity every vertex in $W^-$ to $t$.
\item If $D[\{y_i,z_i~|~i \in I_B\}]$ is transitively equivalent to $O_1^{\ell_5(t)}$, further identity every vertex in $Y^-$ to $t$.
\end{enumerate}
\end{identification rule}

\begin{identification rule}
\label{ir:WX}
For every consecutive $i < i' \in I_A^2 \cup \{1,t\}$ $($that is, $(I_A^2 \cup \{1,t\}) \cap [i,i'] = \{i,i'\})$ and every $j \in[i,i']$, proceed as follows.
\begin{enumerate}[nosep]
\item $D[\{w_i,x_i~|~i \in I_A\}]$ is transitively equivalent to $O_2^{\ell_5(t)}$: 
if $w_j \in W^\star \cup W^-$ then identify $w_j$ to $w_i$ and if $x_j \in X^+$ then identify $x_j$ to $w_{i'}$.
\item $D[\{w_i,x_i~|~i \in I_A\}]$ is transitively equivalent to $O_3^{\ell_5(t)}$:
if $w_j \in W^-$ then identify $w_j$ to $x_i$ and if $x_j \in X^\star \cup X^+$ then identify $x_j$ to $x_{i'}$.
\item $D[\{w_i,x_i~|~i \in I_A\}]$ is transitively equivalent to $O_4^{\ell_5(t)}$: 
if $w_j \in W^\star \cup W^-$ then identify $w_j$ to $w_i$ and if $x_j \in X^\star \cup X^+$ then identify $x_j$ to $x_{i'}$.
\end{enumerate}
\end{identification rule}

\begin{identification rule}
\label{ir:YZ}
For every consecutive $i < i' \in I_B^2 \cup \{1,t\}$ and every $j \in[i,i']$, proceed as follows.
\begin{enumerate}[nosep]
\item $D[\{y_i,z_i~|~i \in I_B\}]$ is transitively equivalent to $O_2^{\ell_5(t)}$: 
if $y_j \in Y^\star \cup Y^-$ then identify $y_j$ to $y_i$ and if $z_j \in Z^+$ then identify $z_j$ to $y_{i'}$.
\item $D[\{y_i,x_i~|~i \in I_B\}]$ is transitively equivalent to $O_3^{\ell_5(t)}$:
if $y_j \in Y^-$ then identify $y_j$ to $z_i$ and if $z_j \in Z^\star \cup Z^+$ then identify $z_j$ to $z_{i'}$.
\item $D[\{y_i,z_i~|~i \in I_B\}]$ is transitively equivalent to $O_4^{\ell_5(t)}$: 
if $y_j \in Y^\star \cup Y^-$ then identify $y_j$ to $y_i$ and if $z_j \in Z^\star \cup Z^+$ then identify $z_j$ to $z_{i'}$.
\end{enumerate}
\end{identification rule}

Let $D_I$ be the digraph resulting from an exhaustive application of Identification Rules~\ref{ir:source}-\ref{ir:YZ}. 
We aim to show that $D_I$ can be identified to a digraph which is transitively equivalent to an almost $\sizeXYt$-hard matching pattern. To this end, we first prove the following claims and then
distinguish cases depending on the graphs to which $D[\{w_i,x_i~|~i \in I_A\}]$ and $D[\{y_i,z_i~|~i \in I_B\}]$ are transitively equivalent, and whether there is no edge or every edge in $D$ from $\{w_i~|~i \in I_A^2\}$ to $\{z_i~|~i \in I_B^2\}$ (from $\{y_i~|~i \in I_A^2\}$ to $\{x_i~|~i \in I_B^2\}$, respectively).

\begin{claim}
\label{clm:vertexsetDI}
The vertex set of $D_I$ consists of $\{w_i,x_i~|~i \in I_A^2\} \cup \{y_i,z_i~|~i \in I_B^2\}$ together with possibly $s$ or $t$. 
\end{claim}

\begin{claimproof}
We first prove that $V(D_I) \cap W = \{w_i~|~i \in I_A^2\}$. 
Clearly, after \Cref{ir:source}.1 has been applied, there are no more vertices from $W^\circ \cup W^+$ left.
Now if $D[\{w_i,x_i~|~i \in I_A\}]$ is transitively equivalent to $O_1^{\ell_5(t)}$
then $W^\star = \emptyset$ by \Cref{clm:independent} and once \Cref{ir:sink}.2 has been applied, there are no more vertices from $W^-$ left.
Similarly, if $D[\{w_i,x_i~|~i \in I_A\}]$ is transitively equivalent to $O_3^{\ell_5(t)}$
then $W^\star = \emptyset$ by \Cref{clm:independent} and once \Cref{ir:WX}.2 has been exhaustively applied, 
there are no more vertices from $W^-$ left. 
Otherwise, $D[\{w_i,x_i~|~i \in I_A\}]$ is transitively equivalent to either $O_2^{\ell_5(t)}$ or $O_4^{\ell_5(t)}$ in which case, an exhaustive application of \Cref{ir:WX}.1 and \Cref{ir:WX}.3, respectively, leaves no vertex from $W^\star \cup W^-$.
We conclude symmetrically that $V(D_I) \cap Y = \{y_i~|~i \in I_B^2\}$.

Second, let us show that $V(D_I) \cap X = \{x_i~|~i \in I_A^2\}$.
Clearly, after \Cref{ir:sink}.1 has been applied, there are no more vertices from $X^\circ \cup X^-$ left. 
Now if $D[\{w_i,x_i~|~i \in I_A\}]$ is transitively equivalent to $O_1^{\ell_5(t)}$
then $X^\star = \emptyset$ by \Cref{clm:independent} and once \Cref{ir:source}.2 has been applied, there are no more vertices from $X^+$ left.
Similarly, if $D[\{w_i,x_i~|~i \in I_A\}]$ is transitively equivalent to $O_2^{\ell_5(t)}$
then $X^\star = \emptyset$ by \Cref{clm:independent} and once \Cref{ir:WX}.1 has been exhaustively applied, 
there are no more vertices from $X^+$ left. 
Otherwise, $D[\{w_i,x_i~|~i \in I_A\}]$ is transitively equivalent to either $O_3^{\ell_5(t)}$ or $O_4^{\ell_5(t)}$ in which case, an exhaustive application of \Cref{ir:WX}.2 and \Cref{ir:WX}.3, respectively, leaves no vertex from $X^\star \cup X^+$.
We conclude symmetrically that $V(D_I) \cap Z = \{z_i~|~i \in I_B^2\}$.
\end{claimproof}

\begin{claim}
\label{clm:sourcesink}
If $s$ exists then $d^-_{D_I}(s) = 0$ and if $t$ exists then $d^+_{D_I}(t) = 0$.
\end{claim}

\begin{claimproof}
Assume that $s$ exists and suppose to the contrary that there exists $u \in V(D_I)$ such that $u \in N^-_{D_I}(s)$.
Then there exist $v_s,v_u \in V(D)$ such that $v_s$ has been identified to $s$, $v_u$ has been identified to $u$ and $(v_u,v_s) \in E(D)$. 
Since the source vertex $\sss$ of $D$ (if any) has in-degree zero in $D$, necessarily $v_s \neq \sss$.
Furthermore, $v_s$ cannot belong to $W^\circ \cup W^+ \cup Y^\circ \cup Y^+$: indeed, if, say, $v_s \in W^\circ \cup W^+$ (the case where $v_s \in Y^\circ \cup Y^+$ is symmetric) then by \Cref{clm:W+}, $v_u \in W^\circ \cup W^+$, a contradiction as $v_u$ is not identified to $s$.
Thus, it must be that $D[\{w_i,x_i~|~i \in I_A\}]$ is transitively equivalent to $O_1^{\ell_5(t)}$ and $v_s \in X^+$,
or $D[\{y_i,x_i~|~i \in I_B\}]$ is transitively equivalent to $O_1^{\ell_5(t)}$ and $v_s \in Z^+$.
Assume without loss of generality that the latter holds (the other case is symmetric).
Then since $(A,B)$ is an ordered \toughpair, necessarily $v_u \in Y \cup Z$.
However, since $v_u$ is not identified to $s$, $v_u \notin Y^\circ \cup Y^+$ 
and since $D*[\{y_i|~i \in I_B\}]$ has no edge, $Y^\star = \emptyset$ by \Cref{clm:independent}.
Furthermore, $v_u$ cannot belong to $\{y_i~|~i \in I_A^2\} \cup Y^-$ for otherwise, since $v_s \in Z^+$ and $(A,B)$ is an ordered \toughpair, there exist $i,j \in I_A^2$ such that $i < j$ and $(y_i,z_j) \in E(D^*)$, a contradiction to the fact that $D[\{y_i,x_i~|~i \in I_B\}]$ is transitively equivalent to $O_1^{\ell_5(t)}$.
Thus, it must be that $v_u \in Z$.
However, since $v_s \in Z^+$, $v_u \notin Z^\star \cup Z^- \cup Z^\circ$ by \Cref{clm:independent}
and $v_u \notin \{z_i~|~i \in I_B^2\}$ by definition of $Z^+$. 
Thus, we conclude that $v_u \in Z^+$, a contradiction since $v_u$ is not identified to~$s$.

Assume now that $t$ exists and suppose to the contrary that there exists $u \in V(D_I)$ such that $u \in N^+_{D_I}(t)$.
Then there exist $v_t,v_u \in V(D)$ such that $v_t$ is identified to $t$, $v_u$ is identified to $u$ and $(v_t,v_u) \in E(D)$.
Since the sink vertex $\ttt$ of $D$ (if any) has out0degree zero in $D$, necessarily $v_t \neq \ttt$.
Furthermore, $v_t$ cannot belong to $X^\circ \cup X^- \cup Z^\circ \cup Z^-$: indeed, if, say, $v_t \in X^\circ \cup X^-$ (the case where $v_t \in Z^\circ \cup Z^-$ is symmetric) then by \Cref{clm:X-}, $v_u \in X^\circ \cup X^-$,
a contradiction since $v_u$ is not identified to $t$.
Thus, it must be that $D[\{w_i,x_i~|~i \in I_A\}]$ is transitively equivalent to $O_1^{\ell_5(t)}$ and $v_t \in W^-$, or
$D[\{y_i,z_i~|~i \in I_B\}]$ is transitively equivalent to $O_1^{\ell_5(t)}$ and $v_t \in Y^-$.
Assume without loss of generality that the latter holds (the other case is symmetric).
Then since $(A,B)$ is an ordered \toughpair, necessarily $v_u \in Y \cup Z$.
However, since $v_u$ is not identified to $t$, $v_u \notin Z^\circ \cup Z^-$
and since $D*[\{z_i|~i \in I_B\}]$ has no edge, $Z^\star = \emptyset$ by \Cref{clm:independent}.
Furthermore, $v_u$ cannot belong to $\{z_i~|~i \in I_B^2\} \cup Z^+$ for otherwise, since $v_t \in Y^-$ and $(A,B)$ is an ordered \toughpair, there exist $i,j \in I_B^2$ such that $i < j$ and $(y_i,z_j) \in E(D^*)$, a contradiction to the fact that $D[\{y_i,x_i~|~i \in I_B\}]$ is transitively equivalent to $O_1^{\ell_5(t)}$.
Thus, it must be that $v_u \in Y$.
However, since $v_t \in Y^-$, $v_u \notin Y^\star \cup Y^\circ \cup Y^+$ by \Cref{clm:independent}
and $v_u \notin \{y_i~|~i \in I_B^2\}$ by definition of $Y^-$.
Thus, we conclude that $v_u \in Y^-$, a contradiction since $v_u$ is not identified to $t$. 
\end{claimproof}

\begin{claim}
\label{clm:sameinDIWX}
For any $i \in [4]$, if $D[\{w_j,x_j~|~j\in I_A\}]$ is transitively equivalent to $O_i^{\ell_5(t)}$ then
$D_I[\{w_j,x_j~|~j \in I_A^2\}]$ is transitively equivalent to $O_i^{4g(t)+2}$.
\end{claim}

\begin{claimproof}
Let us first show that if $D[\{w_i~|~i \in I_A\}]$ has no edge then $D_I[\{w_i~|~i \in I_A^2\}]$ has no edge as well. 
Assume that $D[\{w_i~|~i \in I_A\}]$ has no edge, that is, $D[\{w_i,x_i~|~i \in I_A\}]$ is transitively equivalent to $O_1^{\ell_5(t)}$ or $O_3^{\ell_5(t)}$.
Then by \Cref{clm:independent}, $W^\star = \emptyset$.
Suppose first that $D[\{w_i,x_i~|~i \in I_A\}]$ is transitively equivalent to $O_1^{\ell_5(t)}$. 
Then every vertex in $W^\circ \cup W^+$ is identified to $s$ and every vertex in $W^-$ is identified to $t$ 
where $d^-_{D_I}(s) = d^+_{D_I}(t) = 0$ by \Cref{clm:sourcesink}.
Since, in this case, every vertex in $X$ is identified either to $s$, $t$ or itself 
and no other identification may create an edge between two vertices of $\{w_i~|~i \in I_A^2\}$,
we conclude that $D_I[\{w_i~|~i \in I_A^2\}]$ has no edge.
Suppose second that $D[\{w_i,x_i~|~i \in I_A\}]$ is transitively equivalent to $O_3^{\ell_5(t)}$.
Then every vertex in $W^\circ \cup W^+$ is identified to $s$ and every vertex in $X^\circ \cup X^-$ is identified to $t$
where $d^-_{D_I}(s) = d^+_{D_I}(t) = 0$ by \Cref{clm:sourcesink}.
Since every vertex in $W^- \cup X^\star \cup X^+$ is identified to a vertex in $\{x_i~|~i \in I_A^2\}$ 
and no other identification may create an edge between two vertices of $\{w_i~|~i \in I_A^2\}$,
we conclude that $D_I[\{w_i~|~i \in I_A^2\}]$ has no edge.

Second, let us show that if $D[\{w_i~|~i \in I_A\}]$ is transitively equivalent to the directed path 
$w_{\min I_A} \rightarrow w_{\min I_A +1} \rightarrow \ldots \rightarrow w_{\max I_A}$ then
$D_I[\{w_i~|~i \in I_A^2\}]$ is transitively equivalent to the directed path 
$w_{\min I_A^2} \rightarrow w_{\min I_A^2 +1} \rightarrow \ldots \rightarrow w_{\max I_A^2}$.
Assume that $D[\{w_i~|~i \in I_A\}]$ is transitively equivalent to this directed path, 
that is, $D[\{w_i,x_i~|~i \in I_A\}]$ is transitively equivalent to $O_2^{\ell_5(t)}$ or $O_4^{\ell_5(t)}$.
Note that since $I_A^2 \subseteq I_A$, it is enough to show that no edge from a vertex in $\{w_i~|~i \in I_A^2\}$ to a vertex in $\{w_i~|~i \in I_A^2\}$ of smaller index is created (we refer to such an edge as a \emph{bad edge} in the following).
First, since every vertex $w_j \in W^\star \cup W^-$ is identified to the vertex $w_{j'} \in \{w_i~|~i \in I_A^2\}$ such that $j' = \max \{i \in I_A^2~:~ i \leq j\}$, no bad edge is created here.
Now if $D[\{w_i,x_i~|~i \in I_A\}]$ is transitively equivalent to $O_2^{\ell_5(t)}$ (note that, in this case, $X^\star = \emptyset$ by \Cref{clm:independent}), then every vertex $x_j \in X^+$ is identified to the vertex $w_{j'} \in \{w_i~|~i \in I_A^2\}$ such that $j' = \min \{i \in I_A^2~:~j \leq i\}$ and so, no bad edge is created here; and if $D[\{w_i,x_i~|~i \in I_A\}]$ is transitively equivalent to $O_4^{\ell_5(t)}$, then every vertex in $X^\star \cup X^+$ is identified to a vertex in $\{x_i~|~i \in I_A^2\}$ and thus, no bad edge is created here as well. 
Since every vertex in $W^\circ \cup W+$ is identified to $s$ and every vertex in $X^\circ \cup X^-$ is identified to $t$ where $d^-_{D_I}(s) = d^+_{D_I}(t) = 0$ by \Cref{clm:sourcesink}, no bad edge is created with these identifications;
and since no other identification may create a bad edge, we conclude that $D_I[\{w_i~|~i \in I_A^2\}]$ is indeed transitively equivalent to the directed path $w_{\min I_A^2} \rightarrow \ldots \rightarrow w_{\max I_A^2}$.

Consider next $X$. As in the case of $W$, let us first show that if $D[\{x_i~|~i \in I_A\}]$ has no edge then $D_I[\{x_i~|~i \in I_A^2\}]$ has no edge as well. 
Assume that $D[\{x_i~|~i \in I_A\}]$ has no edge, that is, $D[\{w_i,x_i~|~i \in I_A\}]$ is transitively equivalent to $O_1^{\ell_5(t)}$ or $O_2^{\ell_5(t)}$.
Then by \Cref{clm:independent}, $X^\star = \emptyset$.
Now if $D[\{w_i,x_i~|~i \in I_A\}]$ is transitively equivalent to $O_1^{\ell_5(t)}$,
we conclude, symmetrically to the case of $W$, that $D_I[\{x_i~|~i \in I_A^2\}]$ has no edge.
Suppose therefore that $D[\{w_i,x_i~|~i \in I_A\}]$ is transitively equivalent to $O_2^{\ell_5(t)}$.
Since then, every vertex in $X^\circ \cup X^-$ is identified to $t$, every vertex in $W^\circ \cup W^+$ is identified to $s$  where $d^-_{D_I}(s) = d^+_{D_I}(t) = 0$ by \Cref{clm:sourcesink},
and every vertex in $X^+ \cup W^\star \cup W^-$ is identified to a vertex in $\{w_i~|~i \in I_A^2\}$, 
no edge between two vertices in $\{x_i~|~i \in I_A^2\}$ is created with these identifications.
Since no other identification may create an edge between two vertices in $\{x_i~|~i \in I_A^2\}$, 
we conclude that $D_I[\{x_i~|~i \in I_A^2\}]$ has no edge in this case as well.

Let us next show that if $D[\{x_i~|~i \in I_A\}]$ is transitively equivalent to the directed path 
$x_{\min I_A} \rightarrow x_{\min I_A +1} \rightarrow \ldots \rightarrow x_{\max I_A}$ then
$D_I[\{x_i~|~i \in I_A^2\}]$ is transitively equivalent to the directed path 
$x_{\min I_A^2} \rightarrow x_{\min I_A^2 +1} \rightarrow \ldots \rightarrow x_{\max I_A^2}$.
Assume that $D[\{x_i~|~i \in I_A\}]$ is transitively equivalent to this directed path, that is,
$D[\{w_i,x_i~|~i \in I_A\}]$ is transitively equivalent to $O_3^{\ell_5(t)}$ or $O_4^{\ell_5(t)}$.
Note that since $I_A^2 \subseteq I_A$, it is enough to show that no edge from a vertex in $\{x_i~|~i \in I_A^2\}$ to a vertex in $\{x_i~|~i \in I_A^2\}$ of smaller index is created (we refer to such an edge as a \emph{bad edge} in the following).
First, since every vertex $x_j \in X^\star \cup X^+$ is identified to the vertex $x_{j'} \in \{w_i~|~i \in I_A^2\}$ such that $j' = \min \{i \in I_A^2~:~ j \leq i\}$, no bad edge is created here. 
Now if $D[\{w_i,x_i~|~i \in I_A\}]$ is transitively equivalent to $O_3^{\ell_5(t)}$ (note that, in this case, $W^\star = \emptyset$ by \Cref{clm:independent}), then every vertex $w_j \in W^-$ is identified to the vertex $x_{j'} \in \{x_i~|~i \in I_A^2\}$ such that $j' = \max \{i \in I_A^2~:~i \leq j\}$ and so, no bad edge is created here; and if $D[\{w_i,x_i~|~i \in I_A\}]$ is transitively equivalent to $O_4^{\ell_5(t)}$, then every vertex in $W^\star \cup W^-$ is identified to a vertex in $\{w_i~|~i \in I_A^2\}$ and thus, no bad edge is created here as well. 
Since every vertex in $X^\circ \cup X^-$ is identified to $t$ and every vertex in $W^\circ \cup W+$ is identified to $s$ where $d^-_{D_I}(s) = d^+_{D_I}(t) = 0$ by \Cref{clm:sourcesink}, no bad edge is created with these identifications;
and since no other identification may create a bad edge, we conclude that $D_I[\{x_i~|~i \in I_A^2\}]$ is indeed transitively equivalent to the directed path $x_{\min I_A^2} \rightarrow \ldots \rightarrow x_{\max I_A^2}$.

By combining the above four properties, we conclude that if $D[\{w_j,x_j~|~j\in I_A\}]$ is transitively equivalent to $O_i^{\ell_5(t)}$ for some $i \in [4]$, then
$D_I[\{w_j,x_j~|~j \in I_A^2\}]$ is transitively equivalent to $O_i^{|I_A^2|}$; and since $|I_A^2| = 4g(t) + 2$, the claims follows.
\end{claimproof}

Symmetrically to \Cref{clm:sameinDIWX}, we have the following.

\begin{claim}
\label{clm:sameinDIYZ}
For any $i \in [4]$, if $D[\{y_j,z_j~|~j\in I_B\}]$ is transitively equivalent to $O_i^{\ell_5(t)}$ then
$D_I[\{y_j,z_j~|~j \in I_B^2\}]$ is transitively equivalent to $O_i^{4g(t)+2}$.
\end{claim}

Now note that since every vertex in $A$ ($B$, respectively) is identified either to $s$, $t$ or a vertex in $A$ ($B$, respectively) and $d^-_{D_I}(s) = d^+_{D_I}(t) = 0$ by \Cref{clm:sourcesink}, 
no edge between $A$ and $B$ is created by these identifications that didn't already exist.
This implies in particular that if there is no edge in $D$ from $\{w_i~|~i \in I_A^2\}$ to $\{z_i~|~i \in I_B^2\}$ 
and from $\{y_i~|~i\in I_B^2\}$ to $\{x_i~|~i \in I_A^2\}$ then by Claims~\ref{clm:sameinDIWX} and \ref{clm:sameinDIYZ}, $D_I - \{s,t\}$ is transitively equivalent to the disjoint union of $O_i^{4g(t)+2}$ and $O_j^{4g(t)+2}$, 
where $O_i^{4g(t)+2}$ ($O_j^{4g(t)+2}$, respectively) is the digraph to which $D[\{w_a,x_a~|~a \in I_A\}]$ ($D[\{y_a,z_a~|~a \in I_B^2\}]$, respectively) is transitively equivalent.
Thus, since we may safely identify, e.g., the two first edges of $D[\{w_i,x_i~|~i \in I_A^2\}]$ 
(that is, $(w_{\min I_A^2},x_{\min I_A^2})$ and $(w_{\min I_A^2+1},x_{\min I_A^2+1})$) to the third edge of
$D[\{w_i,x_i~|~i \in I_A^2\}]$ and proceed similarly with $D[\{y_i,z_i~|~i \in I_B^2\}]$,
we conclude that, in this case, $D_I$ can be identified to an almost $4g(t)$-hard matching pattern.
Note that, more generally, if there is no edge in $D$ from $\{w_i~|~i \in I_A^2\}$ to $\{z_i~|~i \in I_B^2\}$
then $D_I[\{w_i~|~i \in I_A^2\} \cup \{z_i~|~i \in I_B^2\}]$ is transitively equivalent to the disjoint union of $D_I[\{w_i~|~i \in I_A^2\}]$ and $D_I[\{z_i~|~i \in I_B^2\}]$ (in particular, no vertex $r_{WZ}$ may be created by identifications); 
and the same holds for $\{y_i~|~i\in I_B^2\}$ and $\{x_i~|~i \in I_A^2\}$.

Assume henceforth that there is every edge either from $\{w_i~|~i \in I_A^2\}$ to $\{z_i~|~i \in I_B^2\}$,
from $\{y_i~|~i\in I_B^2\}$ to $\{x_i~|~i \in I_A^2\}$, or both.
Note that since we assume $D$ not to contain a $g(t)$-induced-biclique,
if there is every edge from $\{w_i~|~i \in I_A^2\}$ to $\{z_i~|~i \in I_B^2\}$
then at least one of $D_I[\{w_i~|~i \in I_A^2\}]$ and $D_I[\{z_i~|~i \in I_B^2\}]$ has edges 
(in fact, is transitively equivalent to a directed path by Claims~\ref{clm:sameinDIWX} and \ref{clm:sameinDIYZ});
and the same holds for $\{y_i~|~i\in I_B^2\}$ and $\{x_i~|~i \in I_A^2\}$.
Now if there is every edge in $D$ from $\{w_i~|~i \in I_A^2\}$ to $\{z_i~|~i \in I_B^2\}$
and both $D_I[\{w_i~|~i \in I_A^2\}]$ and $D_I[\{z_i~|~i \in I_B^2\}]$ are transitively equivalent to a directed path,
then $D_I\{w_i~|~i \in I_A^2\} \cup \{z_i~|~i \in I_B^2\}]$ is transitively equivalent to the directed path
$w_{\min I_A^2} \rightarrow \ldots \rightarrow w_{\max I_A^2} \rightarrow z_{\min I_B^2} \rightarrow \ldots \rightarrow z_{\max I_B^2}$ in which case we need not introduce the vertex $r_{WZ}$;
and we conclude symmetrically with $\{y_i~|~i\in I_B^2\}$ and $\{x_i~|~i \in I_A^2\}$.
 
By the above discussion, there remains to consider the case where every edge in $D$ from $\{w_i~|~i \in I_A^2\}$ to $\{z_i~|~i \in I_B^2\}$ (or from $\{y_i~|~i\in I_B^2\}$ to $\{x_i~|~i \in I_A^2\}$) exist and exactly one of $D_I[\{w_i~|~i \in I_A^2\}]$ and $D_I[\{z_i~|~i \in I_B^2\}]$ ($D_I[\{y_i~|~i\in I_B^2\}]$ and $D_I[\{x_i~|~i \in I_A^2\}]$, respectively) has no edge. In this case, we devise the following identification rules (for simplicity, we still call $D_I$ the graph resulting from these identifications).

\begin{identification rule}
\label{rr:WZ}
If there is every edge from $\{w_i~|~i \in I_A^2\}$ to $\{z_i~|~i \in I_B^2\}$ in $D_I$ 
and exactly one of $D_I^*[\{w_i~|~i \in I_A^2\}]$ and $D_I^*[\{z_i~|~i \in I_B^2\}]$ has no edge, then proceed as follows.
\begin{enumerate}[nosep]
\item If $D_I^*[\{w_i~|~i \in I_A^2\}]$ has no edge then identify $y_{\min I_B^2}$ to $s$ and
rename $z_{\min I_B^2}$ as $r_{WZ}$.
\item If $D_I^*[\{z_i~|~i \in I_B^2\}]$ has no edge then identify $x_{\max I_A^2}$ to $t$ and
rename $w_{\max I_A^2}$ as $r_{WZ}$.
\end{enumerate}
\end{identification rule}

\begin{claim}
\label{clm:renamingWZ}
If \Cref{rr:WZ} has been applied then it still holds that $d^-_{D_I}(s) = d^+_{D_I}(t) = 0$.
Furthermore, $N^-_{D_I}(r_{WZ}) \setminus \{s\} = \{w_i~|~i \in I_A^2\}$ and $N^+_{D_I}(r_{WZ}) \setminus \{t\} = \{z_i~|~i \in I_B^2 \setminus \{\min I_B^2\}\}$.
\end{claim}

\begin{claimproof}
Assume that \Cref{rr:WZ} has been applied and that $D_I^*[\{w_i~|~i \in I_A^2\}]$ has no edge (the case where $D_I^*[\{z_i~|~i \in I_B^2\}]$ has no edge is symmetric).
Then, by assumption, $D_I^*[\{z_i~|~i \in I_B^2\}]$ must have edges;
and in fact, by \Cref{clm:sameinDIYZ} and definition of $I_B^2$, 
$D_I^*[\{z_i~|~i \in I_B^2\}]$ must be transitively equivalent to the directed path 
$z_{\min I_B^2} \rightarrow z_{\min I_B^2 +1} \rightarrow \ldots \rightarrow z_{\max I_B^2}$.
It follows that $N^+_{D_I}(r_{WZ}) \setminus \{t\} = N^+_{D_I}(z_{\min I_B^2}) \setminus \{t\} = \{z_i~|~i \in I_B^2 \setminus \{\min I_B^2\}\}$
and since, by assumption, there is every from $\{w_i~|~i \in I_A^2\}$ to $\{z_i~|~i \in I_B^2\}$ in $D_I$,
we conclude that $N^-_{D_I}(r_{WZ}) \setminus \{s\} = \{w_i~|~i \in I_A^2\}$.
Finally to see that $d^-_{D_I}(s) = d^+_{D_I}(t) = 0$, note that by construction, the only possible in-neighbor for $y_{\min I_B^2}$ is $s$ and the only possible out-neighbor for $x_{\max I_A^2}$ is $t$.
\end{claimproof}

\begin{identification rule}
\label{rr:YX}
If there is every edge from $\{y_i~|~i \in I_B^2\}$ to $\{x_i~|~i \in I_A^2\}$ in $D_I$ 
and exactly one of $D_I^*[\{y_i~|~i \in I_B^2\}]$ and $D_I^*[\{x_i~|~i \in I_A^2\}]$ has no edge, then proceed as follows.
\begin{enumerate}
\item If $D_I^*[\{y_i~|~i \in I_B^2\}]$ has no edge then identify $w_{\min I_A^2}$ to $s$ and
rename $x_{\min I_A^2}$ as $r_{YX}$.
\item If $D_I^*[\{x_i~|~i \in I_A^2\}]$ has no edge then identify $z_{\max I_B^2}$ to $t$ and 
rename $y_{\max I_B^2}$ as $r_{YX}$.
\end{enumerate}
\end{identification rule}

Symmetrically to \Cref{clm:renamingWZ}, we have the following.

\begin{claim}
\label{clm:renamingYX}
If \Cref{rr:YX} has been applied then it still holds that $d^-_{D_I}(s) = d^+_{D_I}(t) = 0$.
Furthermore, $N^-_{D_I}(r_{YX}) \setminus \{s\} = \{y_i~|~i \in I_B^2 \setminus \{\max I_B^2\}\}$ and $N^+_{D_I}(r_{YX}) \setminus \{t\} = \{x_i~|~i \in I_A^2\}$.
\end{claim}

Note finally that \Cref{rr:WZ} reduces the size of $E(D_I - \{s,t\})$ by one as does \Cref{rr:YX} (it may be that the size of $\{(w_i,x_i)~|~i \in I_A^2\}$ reduces by two while the size of $\{(y_i,z_i)~|~i \in I_B^2\}$ remains the same if \Cref{rr:WZ}.2 and \Cref{rr:YX}.1 are applied, and vice-versa). Since we may always safely identify consecutive edges in $\{(w_i,x_i)~|~i \in I_A^2\}$ or $\{(y_i,z_i)~|~i \in I_B^2\}$ to reduce their size to $4g(t)$ (as done above), we conclude that $D_I$ can be identified to a digraph transitively equivalent to an almost $4g(t)$-hard matching pattern.\\

There remains to show how to obtain a $g(t)$-hard matching pattern from an almost $4g(t)$-hard matching pattern. To this end, let $H$ be an almost $\sizeXYt$-hard matching pattern. We first show how to handle the source vertex from item 3 in \Cref{def:almost} and then how to handle the sink vertex from item 4 in \Cref{def:almost}.

Let $s$ be the vertex of $H$ such that $N^-(s) = \emptyset$ and $N^+(s) \subseteq W \cup X \cup Y \cup Z$.
Suppose first that $N^+(s) \cap W \neq \emptyset$ and let $I_1 \subseteq [\sizeXYt]$ be the set of indices such that $N^+(s) \cap W = \{w_i~|~i \in I_1\}$. By the pigeonhole principle, either (1) $|I_1| \geq \sizeXYt/2$ or (2) $|[\sizeXYt] \setminus I_1| \geq \sizeXYt/2$. If (1) holds then we proceed as follows.
\begin{itemize}
\item For every $j < \min I_1$, we identify $w_j$ and $x_j$ to $w_{\min I_1}$ and $x_{\min I_1}$, respectively.
\item For every consecutive $i < i' \in I_1$ and for every $j \in ]i,i'[$, we identify $w_j$ and $x_j$ with $w_i$ and $x_i$, respectively.
\item For every $j > \max I_1$, we identify $w_j$ and $x_j$ to $w_{\max I_1}$ and $x_{\max I_1}$, respectively. 
\end{itemize}
Suppose next that (2) holds. If $H$ contains the path $w_1 \rightarrow w_2 \rightarrow \ldots \rightarrow w_{\sizeXYt}$ then we simply identify $s$ to $w_1$. Let us therefore assume that $H$ does not contain this path. If $H$ contains the path $x_1 \rightarrow x_2 \rightarrow \ldots \rightarrow x_{\sizeXYt}$ then for every $i \in I_1$, we identify $w_i$ to $x_i$ and further identify $s$ to $x_1$.
Finally, if $H$ does not contain this path then we identify every vertex in $\{w_i,x_i~|~i \in I_1\} \cup \{s\}$ to $w_1$.

Now suppose that $N^+(s) \cap W = \emptyset$ and $N^+(s) \cap X \neq \emptyset$. Let $I_1 \subseteq [\sizeXYt]$ be the set of indices such that $N^+(s) \cap W = \{w_i~|~i \in I_1\}$. By the pigeonhole principle, either (1) $|I_1| \geq \sizeXYt/2$ or (2) $|[\sizeXYt] \setminus I_1| \geq \sizeXYt/2$. If (1) holds then we proceed as follows.
\begin{itemize}
\item For every $j < \min I_1$, we identify $w_j$ and $x_j$ to $w_{\min I_1}$ and $x_{\min I_1}$, respectively.
\item For every consecutive $i < i' \in I_1$ (that is, $I_1 \cap [i,'i]= \{i,i'\}$) and for every $j \in ]i,i'[$, we identify $w_j$ and $x_j$ with $w_i$ and $x_i$, respectively.
\item For every $j > \max I_1$, we identify $w_j$ and $x_j$ to $w_{\max I_1}$ and $x_{\max I_1}$, respectively. 
\end{itemize}
Suppose next that (2) holds. If $H$ contains the path $x_1 \rightarrow x_2 \rightarrow \ldots \rightarrow x_{\sizeXYt}$ then we simply identify $s$ to $x_1$. Let us therefore assume that $H$ does not contain this path. If $H$ contains the path $w_1 \rightarrow w_2 \rightarrow \ldots \rightarrow w_{\sizeXYt}$ then for every $i \in I_1$, we identify $x_i$ to $w_i$ and further identify $s$ to $w_1$. Finally, if $H$ does not contain this path then we identify every vertex in $\{w_i,x_i~|~i \in I_1\} \cup \{s\}$ to $w_1$.

By proceeding symmetrically with $Y$ and $Z$, we may identify $H$ to an almost $2g(t)$-hard matching pattern $H_1$ in which $s$ satisfies item 3 of \Cref{def:cleanedtoughpair}. 
Now let $t$ be the vertex of $H$ (and of $H_1$) such that $N^+(t) = \emptyset$ and $N^-(t) \subseteq W \cup X \cup Y \cup Z$. Then using similar arguments, we may further identify $H_1$ to an almost $g(t)$-hard matching pattern $H_2$ in which $s$ still satisfies item 3 of \Cref{def:cleanedtoughpair} and $t$ satisfies item 4 of \Cref{def:cleanedtoughpair}, that is, $H_2$ is a $g(t)$-hard matching pattern which concludes the proof. 
\end{proof}